\section{Free monoids and the interval diagram}
\label{sec:G}
\providecommand{\hocolim}{\operatorname{hocolim}}
\providecommand{\gp}{\mathrm{gp}}
\providecommand{\ab}{\mathrm{ab}}
\providecommand{\Und}{\operatorname{Und}}

We keep the distinction between the colimit over all embeddings and the
filtered colimit over inclusions.  In this section $A=\Cdelta$, $M_0$ is
its ordered strong Fra\"iss\'e limit, $G=\Aut(M_0)$, and $I$ is the
directed poset of nonempty finite strong subsets of $M_0$.  The letter
$M$ denotes a diagram of monoids, and we write
\[
 S=\colim_{B\in I}^{\Mon}M(B),\qquad
 C=\colim_{B\in A}^{\Mon}M(B).
\]
The group acts smoothly on $S$ by monoid automorphisms.  Thus $C$ is
the coequalizer of that action in monoids, denoted $S_G$, rather than
the set of orbits unless the orbit relation happens to be a congruence.

\subsection{Checks on the inherited reductions}

\begin{proposition}[Known, rechecked: the exact-index comparison]
\label{G:exact}
Let $J$ be a nonempty connected small category with exact integral
colimits.  For every $M:J\to\FreeMon$, the canonical functor
\[
 q_M:\int_J M\longrightarrow
 \mathbf B\bigl(\colim_J^{\Mon}M\bigr)
\]
has a weakly contractible target-under-comma category, and is a weak
equivalence.  This is DF D.6, with arbitrary ranks and arbitrary
unit-preserving substitutions.
\end{proposition}
\begin{proof}
Put $C=\colim_JM$.  The comma is the Grothendieck construction of
the diagram $D(j)$ with object set $C$ and arrows
\[
 c\longrightarrow j_j(m)c\quad(m\in M(j)).
\]
For a basis $E_j$ of $M(j)$ this is the free category on the edges
$c\to j_j(e)c$, $e\in E_j$.  Its integral homology vanishes above
degree one.  This assertion follows directly by removing the first
edge of each nonempty path in the representable resolution; it is
valid for infinite graphs.

The ordinary categorical colimit of the $D(j)$ is the regular
translation category $C//C$.  To check this, present $C$ by all stage
basis letters and all substitution relations.  At each initial
vertex $c$, the same relations are precisely path relations; a word
forces every intermediate vertex by multiplication.  Consequently
there are exactly the arrows labelled $u$ with $uc=c'$.  This
category has initial object $1$.

Exactness implies totality, so $BJ$ is simply connected.  In the
groupoid completion of $\int_JD$, the transport arrows at each
$c\in C$ therefore form a connected thin groupoid.  Choosing their
unique classes eliminates all transports.  The remaining generators
and relations are the groupoid completion of $C//C$, which is
terminal.  The comma is thus connected and simply connected.
Finally,
\[
 E^2_{p,q}=L_p\colim_J H_q(BD(j);\Z)
 \Longrightarrow H_{p+q}(B\!\int_JD;\Z)
\]
has no terms with $p>0$, by exactness, and no terms with $q>1$,
by freeness.  All positive homology of the comma vanishes.  Hurewicz
and Whitehead make it weakly contractible, and Theorem A applies.
This checks the precise place where DF D.6 uses exactness: the
positive-column terms for the cycle-module row.  We have not
replaced that use by totality.
\end{proof}

\begin{proposition}[Known, rechecked: saturation and finite free charts]
\label{G:saturation}
For $M:A\to\Mon$, put $T_S(B)=S^{G_B}$.  The natural monoid map
$M\to T_S$ induces a weak equivalence on the homotopy colimits of
classifying spaces and an isomorphism on ordinary monoid colimits.
For $M:A\to\FreeMon$, every finite subset of $S$ lies in a finite
free submonoid.  It has a least such free hull.  The hull of a finite
set fixed pointwise by a subgroup is itself fixed pointwise.
Consequently the universal comparison question on $A$ reduces to
finite-rank free monoids with injective transitions, as in DF D.8--D.9.
\end{proposition}
\begin{proof}
We record the checks rather than infer these conclusions merely
from a filtered colimit.  Eventual equality in a set diagram is an
equivalence relation by amalgamation.  A finite collection of
equality witnesses is synchronized by successive amalgamation, so
injectivization commutes with finite products.  Every embedding
extends to $G$; applying its inverse shows that eventual equality
is exactly equality in the filtered colimit over inclusions.

For a set action $H$, the comparison to its injectivization is
homotopy final.  After choosing a representative, its comparison
commas reduce, by Theorem A, to nonempty full cosieves in slices
$B/A$.  For the comparison of an injective action with $T_S$, the
comma itself is the cosieve of arrows on whose target the given
element is represented.  It is nonempty because two finite strong
supports lie in a common finite strong set.  Each such cosieve is
contracted by the functorial free amalgam with one fixed object over
$B$: in each order gap put the variable side first.  Strongness of
the induced arrows follows by viewing the target as the free
amalgam of the larger variable side and the source over the smaller
side.  This verifies both Theorem A arguments of T A.33 for the
present use, including their variance.

Apply these comparisons to $H=\Und(M)^k$ for every finite $k$.
Filtered colimits and fixed subsets commute with finite products.
Rowwise weak equivalences in the bisimplicial bar construction give
the homotopy-colimit assertion.  For ordinary colimits, the
isomorphisms of colimits of underlying sets for $k=1,2$ identify
the presentation whose generators are stage elements and whose
relations are transition and multiplication relations.  Hence the
monoid-colimit map is an isomorphism.  This argument does not claim
that arbitrary set colimits commute with products.

Here are the word-theoretic checks of DF D.9.  The filtered colimit
$S$ is cancellative and conical, since any finite equality is
witnessed in one free stage.  Given a finite set $X\subset S$, choose,
over all later stages, a free submonoid containing its image of
least possible rank.  A noninjective subsequent map on that chart
would give a non-code among its basis images; the defect theorem
\cite[Section 2, Theorems 1--2]{GCastiglione2019} puts those images
in a free monoid of smaller rank, a contradiction.
Thus this chart embeds in $S$.

If $H,L$ are free submonoids of $S$, then $H\cap L$ is stable
inside $H$.  Indeed if $p,up,q,qu\in H\cap L$, apply Levi's lemma
in $L$ to $q(up)=(qu)p$.  One factorization gives an element
$t\in L$ with $t=u$ by cancellation; the other gives $ut=1$,
and conicality gives $u=1$.  Hence $u\in H\cap L$.  The stability
criterion, in the opened source \cite[Theorem 1, p.~84]{GCarpi2014},
makes this intersection free.  It follows that the ordinary free
hull of $X$ in one chart lies in every free submonoid of $S$
containing $X$.  It is therefore the intrinsic least free hull.
Every atom of the hull occurs in a word of $X$; an automorphism
fixing $X$ fixes those ordered letters, and hence the hull pointwise.

Replace $M(B)$ by the free hull in $S$ of its finite basis image.
The hulls are functorial under embeddings extended to $G$, all
maps are injective, and their filtered union is $S$.  Saturation
therefore identifies both comparisons.  Finally an arbitrary-rank
diagram is the filtered union of subdiagrams generated by finitely
many tagged letters and the basis supports of all their images.
At any object these are finite because the hom-sets of $A$ are
finite.  Positivity of word substitutions makes them subdiagrams.
Nerves, monoid colimits, and weak equivalences of simplicial sets
commute with these filtered unions.  This proves the reduction.
\end{proof}

\begin{proposition}[Proved: the orbit-divisibility obstruction]
\label{G:divisibility}
Let $M:A\to\FreeMon$.  Suppose that in its smooth filtered-limit
monoid $S$ no element is a proper left multiple of another element
in its $G$-orbit:
\[
 x=(g x)w\quad\Longrightarrow\quad w=1
 \qquad(x,w\in S, g\in G).
 \tag{G.1}
\]
Then $q_M$ is a weak equivalence and
$L_i\colim_A(M_{\ab})=0$ for every $i>0$.
In particular failure of either conclusion forces a relation
$x=(g x)w$ with $w\ne1$.  A strictly positive invariant additive
length implies (G.1), but is not needed in this proposition.
\end{proposition}
\begin{proof}
Finite free charts make $S$ equidivisible.  If $g(ab)=ab$, Levi's
lemma makes $ga$ and $a$ comparable by left divisibility.  Condition
(G.1), also applied with $g^{-1}$, forces $ga=a$, and cancellation
forces $gb=b$.  Therefore every fixed monoid $S^{G_B}$ is closed
under taking factors in $S$.

We check the rest of DF D.10 with exactly this hypothesis.  Let $J$
be the directed poset of finite free submonoids $H\subset S$.
Directedness follows from the finite-chart assertion above.
Factor-closedness makes $H\cap S^{G_B}$ the free submonoid on
the subset of atoms of $H$ fixed by $G_B$.
Fix a monoid $C_0$ with trivial $G$-action and a $G$-invariant
homomorphism $\rho:S\to C_0$.  The free translation category
of $H\cap S^{G_B}$ on the set $C_0$ has graph-chain group
\[
 E_{H,B}=\Z[C_0\times(\operatorname{At}(H)\cap S^{G_B})]
\]
in degree one and $\Z[C_0]$ in degree zero.  An inclusion of charts
sends each edge to the chain of its substituted path.  These maps
compose strictly.  Their filtered colimit $E_B$ is a submodule of
the corresponding full group $E$, since filtered colimits are exact.

At each chart the subgroups in this construction are spanned by
subsets of one basis.  The identity
$\langle G_B,G_D\rangle=G_{B\cap D}$ therefore gives
\[
 E_B\cap\sum_jE_{D_j}=\sum_jE_{B\cap D_j}.
 \tag{G.2}
\]
The passage to the filtered colimit preserves these finite
intersections.  Filtering $E_B$ by the cardinality of a strong
support gives successive quotients
\[
 \bigoplus_{[D],\ |D|=n}\Z[A(D,B)]\otimes V_D
 \quad(n>0),
 \qquad V_\varnothing\quad(n=0).
\]
Here $V_D=E_D/\sum_{E\subsetneq D,\ E\text{ strong}}E_E$;
(G.2) proves both independence of equal-size supports and that
these are all relations.  Rigidity removes a setwise-stabilizer
quotient.  The representable summands are acyclic for higher
colimits for arbitrary $V_D$, using the initial object of $D/A$;
the constant summand is acyclic because $BA$ is contractible.
The exhaustive filtration and exact filtered colimits give
$L_i\colim_A E_-=0$ for $i>0$.

The translation-category graph complexes now give an exact sequence
\[
 0\longrightarrow H_1\longrightarrow E_-
 \longrightarrow\Z[C_0]
 \longrightarrow\Z[\pi_0]\longrightarrow0.
\]
The last term is the free module on an action of $A$, so it is
acyclic by totality.  The constant term and $E_-$ are acyclic as
well.  Splitting this sequence at its middle image gives
$L_i\colim_A H_1=0$ for $i>0$.  Each translation category has no
homology above one, being a filtered colimit of free categories.
Thus its homotopy colimit has no homology above one.

Take $C_0=S_G$.  The same presentation and transport argument as in
Proposition~\ref{G:exact} makes the comparison comma connected and
simply connected.  Its higher homology has just vanished, so it is
weakly contractible.  Saturation transfers this to $q_M$.
Taking instead $C_0=1$ identifies $E_B$ with the universal abelian
group of $S^{G_B}$.  The identical calculation, followed by the
two-row homology spectral sequence for free monoids, gives the
asserted vanishing for $M_{\ab}$.
This supplies the additional argument needed for the necessary
relation stated in the prompt: it is a consequence of the
factor-closedness proof, not of the length statement by itself.
\end{proof}

\begin{proposition}[Proved: rank at most one cannot fail]
\label{G:rankone}
If every $M(B)$ has free rank at most one, then $q_M$ is a weak
equivalence and $L_i\colim_A M_{\ab}=0$ for $i>0$.  Erasing
transitions and rank-zero values are allowed.
\end{proposition}
\begin{proof}
After injectivization, a value is either trivial or the free monoid
on one letter: if a positive power is eventually killed, the whole
generator is eventually killed in a free target; otherwise every
transition from that value is a positive power map.  The filtered
limit is either trivial or is isomorphic to an additive submonoid
of $\Q_{\ge0}$ whose group completion is a nonzero subgroup of
$\Q$.  This follows by fixing one nontrivial stage and writing the
positive power maps between its later stages as rational changes
of scale.

Every automorphism of that monoid acts on its group completion by
one positive rational scale.  The action of $G$ has open kernel,
because fixing any nonzero element forces that scale to be one.
There is no proper open normal subgroup of $G$.  Indeed (Op)
writes it as $G_F$; normality and $\Fix_{M_0}(G_F)=F$ make $F$
a finite invariant set.  Free disjoint copies show that every
point has infinite orbit, so $F=\varnothing$.  Hence $G$ acts
trivially on $S$.

Its fixed-point diagram is now constant with value $S$.  Its
homotopy colimit is $BA\times BS$, and its monoid colimit is
$S$.  The projection is a weak equivalence since $BA$ is
contractible.  Saturation proves the claim about $q_M$.  Finally
$BS$ is a filtered colimit of classifying spaces of monogenic
free monoids, so it has zero homology above degree one; the
free-monoid homology formula gives the derived-colimit assertion.
\end{proof}

\subsection{The gap diagram, including all raw-point types}

\begin{theorem}[Proved: the specified gap comparison passes]
\label{G:gap}
For $B\in A$, let $M_{\rm gap}(B)$ be the free monoid on the gaps
of $\{-\infty\}\cup B\cup\{+\infty\}$, and let an increasing
strong embedding send a gap to the product of all intervening gaps.
Then its specified comparison $q_{M_{\rm gap}}$ is a weak equivalence.
Let $O=G\backslash M_0$ be the set of orbits of \emph{all raw points}.
There are natural identifications
\[
 (M_{\rm gap})_{\ab}\cong\Z[\Und(B)]\oplus\underline\Z,
 \qquad L_i\colim_A(M_{\rm gap})_{\ab}=0\quad(i>0),
\]
and both sides of its comparison have type
$K(F(O\amalg\{*\}),1)$.
In the full category $O$ need not be a singleton; replacing raw
points by $A(V,-)$ would be incorrect.
\end{theorem}
\begin{proof}
Write the gaps of $b_1<\cdots<b_n$ as $e_0,\ldots,e_n$.
In the group completion put
\[
 p_{b_i}=e_0\cdots e_{i-1},\qquad t=e_0\cdots e_n.
\]
They are a free basis: the inverse expressions are
$e_0=p_{b_1}$, $e_i=p_{b_i}^{-1}p_{b_{i+1}}$ for $0<i<n$,
and $e_n=p_{b_n}^{-1}t$.  Subdivision sends $p_b$ to $p_{f(b)}$
and $t$ to $t$.  Thus the group-completed diagram is naturally
the free-group diagram on the raw-point action together with a
constant singleton.  Abelianization gives the displayed natural
direct sum, and totality makes all its positive derived colimits
zero.  More strongly the homotopy colimit of a basis-map free-group
diagram over a total category is the bouquet on the colimit of its
bases.  To see this without a homology-only argument, use the
natural graph model: it is the homotopy pushout obtained by attaching
$B\El(L)\times[0,1]$ to $BA$ at both ends.  Contract each action
component and $BA$.  The resulting bouquet is the specified map
to the free group on $\colim L$.  Group completion of each free
monoid is a weak equivalence, naturally.  Hence
\[
 B\!\int_A M_{\rm gap}\simeq K(F(O\amalg\{*\}),1).
 \tag{G.3}
\]

We must still identify the ordinary \emph{monoid} colimit.
The filtered monoid $S$ has generators $[x,y]$ for
$x<y$ in $M_0\cup\{-\infty,+\infty\}$, with relations
\[
 [x,y][y,z]=[x,z]\quad(x<y<z).
 \tag{G.4}
\]
Every finite set of endpoints lies in a finite strong set, and
subdivision in that set proves both the presentation and its
independence of the chosen strong set.  Every embedding extends
to $G$, so the colimit over all embeddings is the action
coequalizer $S_G$.

Let $\kappa\in O$ be the orbit of a strong singleton.  Free
amalgamation inserts any finite collection of new isolated
$\kappa$-points, in any specified order gaps, over a given finite
strong set.  These points and every subset of them are strong.
For $a\in M_0$, its finite strong closure $\operatorname{cl}(a)$
is the intersection of the strong subsets of one finite strong
set containing it.  Its definition is independent of that set,
by the intersection property.  An automorphism fixing $a$
preserves this finite ordered closure and hence fixes it pointwise.
Consequently the orbit $\alpha=[a]\in O$ determines the marked
ordered closure $(\operatorname{cl}(a),a)$.

Choose independent isolated points $z<w$.  In $S_G$ denote
$[-\infty,z]$ by $l$, $[z,+\infty]$ by $r$, and $[z,w]$ by
$e$.  Their values are independent of the chosen copies, by
homogeneity.  For $\alpha=[a]$, choose new isolated points
$z<a<w$, with $z$ and $w$ in the immediate gaps below and above
$a$ in its closure, and put
\[
 L_\alpha=[z,a],\qquad R_\alpha=[a,w].
\]
These also depend only on $\alpha$.  Subdividing with one or two
additional isolated points gives
\[
\begin{gathered}
 e^2=e,\qquad le=l,\qquad er=r,\\
 eL_\alpha=L_\alpha,\qquad R_\alpha e=R_\alpha,
 \qquad L_\alpha R_\alpha=e,\\
 L_\kappa=R_\kappa=e.
\end{gathered}
\tag{G.5}
\]
For the relation $L_\alpha R_\alpha=e$, the two isolated points
form a strong edgeless pair even though $a$ lies between them.

For arbitrary $x<y$, choose new isolated $c,d$ with
$x<c<d<y$, so close to the two endpoints that $c$ is in the
immediate right gap of $x$ in its closure, and $d$ in the
immediate left gap of $y$ in its closure.  One first works over
a finite strong set containing both closures and then freely
adjoins $c,d$; this proves existence even when the closures
overlap.  In $S_G$ we obtain
\[
 [x,y]=R_{[x]}eL_{[y]}=R_{[x]}L_{[y]}.
\]
The ray and total formulas are
\[
 [-\infty,x]=lL_{[x]},\quad
 [x,+\infty]=R_{[x]}r,\quad
 [-\infty,+\infty]=lr.
 \tag{G.6}
\]
These prove that (G.5) is a presentation, not only a list of
necessary relations.  Namely assign the formulas (G.6) and
$R_{[x]}L_{[y]}$ to the interval generators in an abstract monoid
presented by (G.5).  Each relation (G.4) holds because
$L_\alpha R_\alpha=e$, $R_\alpha e=R_\alpha$, and
$eL_\alpha=L_\alpha$.  This gives a $G$-invariant homomorphism
$S$ to the presented monoid.  The reverse homomorphism sends the
listed generators to their chosen intervals.  The preceding
formulas make the composites identities on every generator.

Let $W_O$ be the category of nonempty finite linearly ordered
$O$-coloured sets and colour-preserving increasing injections.
It is Ramsey, even when $O$ is infinite.  For fixed source and
target only finitely many colours occur.  Number those colours
$1,\ldots,m$ and use the target word $(1,\ldots,m)^N$.
An increasing injection $f:[n]\to[N]$ lifts a word with letters
$c_j$ to positions $m(f(j)-1)+c_j$.  Apply ordinary finite Ramsey
to the colouring of lifted source injections.  The monochromatic
copy of $[n]$ lifts to the required copy of the target word;
every source embedding into that target is one of the lifted
injections.  Thus the morphism Ramsey property holds.  The
Ramsey-to-integral-affine-equalizer theorem R5 makes $W_O$ exact.
Equivalently its explicit affine proof uses the same formula for
lifting the integral cones of $\Delta_{\rm inj}$; each composite
retains the source colour in the final summand, so all affine
identities are preserved.

Colour a raw point $a$ by $[a]\in O$.  This defines a functor
$u:A\to W_O$; the gap diagram on $A$ is its pullback.  The
ordinary colimit of the coloured gap diagram has exactly
presentation (G.5), by the same insertion argument in the dense
coloured order.  Consequently the map on monoid colimits induced
by $u$ is an isomorphism.  The point-action colimit map is the
identity on $O$.  The prefix bases therefore identify the induced
map on the two homotopy colimits with the identity map on the
bouquet of (G.3).  Finally the coloured comparison is a weak
equivalence by Proposition~\ref{G:exact}.  The commutative square
of specified comparisons and two-out-of-three prove the theorem.

The finite certificate \texttt{code/G\_gap\_checks.py} checks
all overlaps of the three-colour presentation, all endpoint-type
subdivision identities, 54,894 colour-preserving injections
through target size six, and 2,295 prefix-basis naturality
equations through eight points.  The unbounded statements are
proved above; the program is a check of the displayed finite
presentation and formulas.
\end{proof}

\subsection{A further necessary obstruction in the target}

\begin{proposition}[Proved: idempotent-free colimits pass]\label{G:idempotent}
Let $M:A\to\FreeMon$, where $A=\Cdelta$, and put
$C=\colim_A^{\Mon}M$.  If $C$ has no idempotent other than $1$,
then its specified comparison $q_M$ is a weak equivalence and
$T_i(M_{\ab})=0$ for every $i>0$.  In particular this holds if
$C$ is left cancellative or right cancellative.

More precisely, use the smooth locally free $G$-monoid $S$ of
Proposition~\ref{G:saturation}, so that $C=S_G$.  If
$x=(gx)w$ with $w\ne1$, then $C$ contains an idempotent $e\ne1$
such that
\[
 [x]e=[x],\qquad e[w]=[w].
 \tag{G.7}
\]
\end{proposition}
\begin{proof}
First, the monoid congruence defining $S_G$ never identifies a
nonidentity element with $1$.  The monoid $S$ is conical: a
product equals $1$ only when both factors are $1$, as is seen in
any free chart containing them.  Each elementary congruence
step has the form
$u s v\leftrightarrow u h(s)v$, for $h\in G$.
One side is $1$ if and only if $u=s=v=1$, and then the other
side is also $1$.  A finite chain of these steps proves the
assertion.

Write $x_0=x$ and $x_1=gx$.  By (Op), write
$\operatorname{Stab}_G(x_0)=G_B$ and
$\operatorname{Stab}_G(x_1)=G_{B'}$, with finite strong supports
$B$ and $B'=gB$.  These supports are unique, by
$\operatorname{Fix}_{M_0}(G_B)=B$.  The support $B$ is nonempty,
since otherwise $gx=x$ and cancellation would give $w=1$.
Moreover $B'\not\subseteq B$.  Indeed they have equal finite
cardinality; if they were equal, $g$ would preserve $B$ setwise,
hence fix it pointwise by the linear order, and would fix $x$.

Choose a finite strong $D$ containing $B$, $B'$ and a support
of $w$.  Freely amalgamate three copies $D_1,D_2,D_3$ of $D$
over $B$.  In each order gap of $B$, put the $D_i\setminus B$
parts in the order $1<2<3$.  Realize this finite strong amalgam
in the limit over $B$.  The copy maps extend to automorphisms
$h_i\in G_B$.  Put
\[
 y_i=h_i x_1,\qquad z_i=h_i w;
 \qquad x_0=y_i z_i\quad(1\le i\le3).
\]
The supports of the $y_i$ are the copies of $B'$.  They are
distinct because $B'\not\subseteq B$ and different copies of
$D$ meet exactly in $B$.  Thus the $y_i$ are distinct.  The
free-chart property gives cancellation and equidivisibility in
$S$, so distinct prefixes of $x_0$ are strictly comparable.

For any $i<j$ and $k<l$, the isomorphism from
$D_i\cup_B D_j$ to $D_k\cup_B D_l$ which matches the ordered
copies and fixes $B$ is an isomorphism of finite strong ordered
structures.  Homogeneity extends it to $G_B$.  Consequently
all pairs $(y_i,y_j)$, $i<j$, are in the same $G_B$-orbit.
Their prefix direction is therefore constant.  If
$y_1<_{\mathrm{pref}}y_2<_{\mathrm{pref}}y_3$, define the
nonidentity quotients by $y_j=y_i p_{ij}$.  Cancellation gives
$p_{13}=p_{12}p_{23}$.  If the direction is reversed, define
$y_i=y_j p_{ij}$ instead, and obtain
$p_{13}=p_{23}p_{12}$.  In either case uniqueness of quotients
and the pair-orbit assertion show that $p_{12},p_{23},p_{13}$
are in a single $G$-orbit.  Their common image $e$ in $C$
satisfies $e^2=e$ and is not $1$, by the first paragraph.

Every $y_i$ is in the $G$-orbit of $x$.  The defining quotient
equations give $[x]e=[x]$.  In the increasing case cancellation
in $y_i z_i=y_j z_j$ gives $z_i=p_{ij}z_j$; the decreasing
case gives $z_j=p_{ij}z_i$.  All $z_i$ are in the orbit of
$w$, so $e[w]=[w]$ in both cases.  This proves (G.7).

If $C$ has no nonidentity idempotent, no proper left factor
can be in the same $G$-orbit.  Proposition~\ref{G:divisibility}
now proves both conclusions.  A left or right cancellative
monoid has no nonidentity idempotent, by cancellation in
$e^2=e$.
\end{proof}

\begin{proposition}[Proved: bounded ranks pass]\label{G:boundedrank}
Let $M:A\to\FreeMon$, where $A=\Cdelta$, and let $S$ be its
smooth filtered-limit monoid.  If $G$ acts trivially on
$S_{\ab}$, then $q_M$ is a weak equivalence and
$T_i(M_{\ab})=0$ for every $i>0$.
The same conclusions hold if the smooth rational representation
$\Q\otimes S_{\ab}$ is locally finite dimensional, meaning that
the linear span of each $G$-orbit is finite dimensional.
In particular they hold if $S_{\ab}$ has finite rational rank,
and for every diagram whose free ranks have one finite uniform
bound.

Conversely, any failure of either conclusion forces a proper
same-orbit prefix and arbitrarily large linearly independent
finite subsets of one $G$-orbit in $\Q\otimes S_{\ab}$.
\end{proposition}
\begin{proof}
The free charts form a directed system under containment: put
the finitely many basis elements of two charts into a third
free chart.  Abelianization commutes with its filtered colimit.
Thus $S_{\ab}$ is a filtered colimit of free abelian groups and
is torsion-free.  Moreover
\[
 s\ne1\quad\Longrightarrow\quad [s]_{\ab}\ne0.
 \tag{G.8}
\]
Indeed, in any free chart containing $s$ its word has a nonzero
nonnegative exponent vector.  The same is true in every later
free chart; hence its class cannot become zero in the filtered
colimit.  This does not assert that the abelianization map is
injective on all of $S$.

If $x=(gx)w$ and the action on $S_{\ab}$ is trivial, then
$[w]_{\ab}=0$.  By (G.8), $w=1$.  Proposition
\ref{G:divisibility} gives the first assertion.

As checked in the proof of Proposition~\ref{G:rankone}, $G$
has no proper open normal subgroup.  For completeness, if an
open normal subgroup is $G_F$ by (Op), normality makes $F$
$G$-invariant, by $\operatorname{Fix}_{M_0}(G_F)=F$.  Every
nonempty finite strong set has infinitely many disjoint copies
by free amalgamation over the empty set, so $F$ must be empty.
Any smooth finite-dimensional rational representation has an
open kernel: intersect the stabilizers of a finite vector-space
basis.  The kernel is normal, so that representation is trivial.
Apply this to the span of each orbit in a locally finite
dimensional representation.  Torsion-freeness then makes the
action on $S_{\ab}$ trivial.  A uniform rank bound on the
original $M(B)$ bounds the dimension of their filtered rational
abelianization colimit, giving the stated special case.

Here is also a direct geometric certificate for the last
assertion.  If $x=(gx)w$ is proper, carry out the construction
in Proposition~\ref{G:idempotent} with any number $n\ge2$ of
free copies $D_i$ over $B=\operatorname{supp}(x)$.
The resulting $y_i$ are distinct comparable prefixes of $x$.
By (G.8), their abelianizations $v_i$ are pairwise distinct.
Write $\operatorname{Stab}_G(v_1)=G_{F_1}$ using (Op).
Then $F_1\subseteq\operatorname{supp}(y_1)\subseteq D_1$:
the inclusion of stabilizers and the fixed-point formula give
the inclusion of supports.  The support $F_1$ is not contained
in $B$, since otherwise the maps between copies over $B$ would
give equal $v_i$.  The supports $F_i$ of the other $v_i$ are
the corresponding copies of $F_1$.

Suppose $\sum_{i=1}^n c_i v_i=0$ with integers $c_i$, not all
zero, and choose $i$ with $c_i\ne0$.  The union $E$ of the
other copies $D_j$, $j\ne i$, is a finite strong free amalgam.
Its pointwise stabilizer fixes every $v_j$ for $j\ne i$,
hence fixes $c_i v_i$.  Torsion-freeness implies that it fixes
$v_i$.  Thus $G_E\subseteq G_{F_i}$, so $F_i\subseteq E$.
But $F_i\subseteq D_i$ and $D_i\cap E=B$, a contradiction.
The $v_i$ are therefore rationally independent.  They belong
to a single $G$-orbit.  If either conclusion fails, the
contrapositive of Proposition~\ref{G:divisibility} supplies
such a proper same-orbit prefix.
\end{proof}

\subsection{A finite obstruction to the direct positive incidence lift}

\begin{proposition}[Proved: the direct positive incidence lift collapses]
\label{G:positiveincidence}
Suppose $M:A\to\FreeMon$ has injective transitions, and choose
$u,v\in M(V)$.  Write $u_i=M(v_i)u$ and $v_i'=M(v_i)v$ in
$M(W)$.  If
\[
 u_2v_3'=u_1v_1',
 \tag{G.9}
\]
then $u_2=u_1$ and $v_3'=v_1'$.
Thus the abelian incidence relator of R3 cannot be lifted by
this direct positive two-word equation while retaining its
nonzero difference classes.  Likewise the alternative proposed
lift $u_2=t u_1$, $v_1'=t v_3'$ cannot extend injectively from
$V,W$ to $A$ when $t\ne1$.

More generally, if only the original group equation
$u_2u_1^{-1}v_3'(v_1')^{-1}=1$ is assumed, then the vector
$[u_2]-[u_1]$ in $M(W)_{\ab}$ is either zero or has both
positive and negative coordinates in its free abelian basis.
\end{proposition}
\begin{proof}
The obstruction is the ordered hypergraph on $1<\cdots<11$
with six edges, indexed from zero:
\[
 \begin{aligned}
 e_0&=\{1,2,9\},& e_1&=\{2,3,10\},& e_2&=\{3,4,11\},\\
 e_3&=\{1,4,5\},& e_4&=\{5,6,7\},& e_5&=\{7,8,9\}.
 \end{aligned}
 \tag{G.10}
\]
It is an object of $A$, every vertex is strong, and every
displayed edge is strong.  Here is an elementary check of all
subsets, independently of the program.  A vertex of this
hypergraph lies in at most two edges.  The edge-intersection
graph is the theta graph with three paths of lengths $1,3,3$
between $e_0$ and $e_3$.  For a set of $r$ edges, let $s$ be
the number of edges induced on those $r$ vertices of the theta
graph.  Their union has $3r-s$ vertices, hence predimension
$3r-2s$.  For $r=1,2,3,4,5,6$, the maximum possible $s$ is
respectively $0,1,2,4,5,7$.  Thus a single edge has
predimension $3$ and every union of at least two edges has
predimension at least $4$.  Adding vertices outside such a
union increases the predimension by $2$ per vertex.  A subset
without an edge has predimension twice its size.  Positivity,
strongness of each singleton, and strongness of each edge
follow.

For each vertex $a$, let $U_a,V_a$ be the word lengths of the
images of $u,v$ in the free monoid $M(X)$, where $X$ is
(G.10).  Equation (G.9) implies, for each $e=(a,b,c)$,
\[
 \tau_e=U_b-U_a=V_a-V_c.
\]
The four $u$-differences on $e_0,e_1,e_2,e_3$ telescope, and
the four $v$-differences on $e_0,e_3,e_4,e_5$ telescope, giving
\[
 \tau_3=\tau_0+\tau_1+\tau_2,\qquad
 \tau_0=\tau_3+\tau_4+\tau_5.
 \tag{G.11}
\]
Consequently $\tau_1+\tau_2+\tau_4+\tau_5=0$.

In the free monoid $M(W)$, (G.9) makes $u_1,u_2$ comparable
as prefixes.  If they differ, write either $u_2=u_1t$ or
$u_1=u_2t$, with $t\ne1$.  Its image along every edge
embedding $W\to X$ is nonempty, by injectivity.  Hence all six
$\tau_e$ are strictly positive in the first case and strictly
negative in the second.  Either conclusion contradicts
(G.11).  Thus $u_1=u_2$, and cancellation gives
$v_1'=v_3'$.  The alternative lift in the statement forces
all six $\tau_e>0$ directly and is excluded by the same
calculation.

For the last assertion, a nonzero vector with nonnegative
coordinates stays nonzero and nonnegative under every edge
map: each basis generator has a nonempty image by injectivity.
Its total coordinate sum in $M(X)_{\ab}$ is therefore strictly
positive for every edge.  A nonpositive vector gives the
opposite sign.  The group equation still gives (G.11), so
either one-sided sign contradicts that identity.

The exact certificate
\texttt{code/G\_positive\_lift\_search.py --explicit} checks
all 2,047 nonempty vertex subsets, all six strong-edge tests,
and the two integral circulation vectors whose sum is
$(0,1,1,0,1,1)$.  Its independent primal linear-program check
also finds the strict length system infeasible.  The integer
identity (G.11), rather than the numerical solver status, is
the proof.  This excludes the specified positive lifts; it does
not exclude every free-monoid realization of a nonexact module.
\end{proof}

\subsection{A tree-order obstruction to positive incidence realizations}

\begin{lemma}[Proved: invariant tree orders on strong singletons]
\label{G:treeorder}
Let $\Omega$ be the $G$-orbit of the strong singleton points of
$M_0$.  A $G$-invariant partial order on $\Omega$ whose principal
downsets are linearly ordered is either equality, the given
linear order, or its reverse.
\end{lemma}
\begin{proof}
Write $\prec$ for its strict relation, and suppose $a\prec b$
with $a\ne b$.  Choose a finite strong $D$ containing $a,b$.
The singletons of both distinguished points are strong in $D$.
Freely amalgamate two copies of $D$ over $b$, putting the first
copy before the second in every gap of $b$.  Call their other
distinguished points $a_1<a_2$.  Both precede $b$ under
$\prec$, so they are comparable.  Denote this two-block
amalgam by $P$.

Now freely amalgamate two copies of $P$, identifying $a_2$ in
the first with $a_1$ in the second.  This is legitimate over
a strong singleton.  The three remaining named positions are
$a_1<a_2<a_3$.  The pairs $(a_1,a_2)$ and $(a_2,a_3)$ are
conjugate, so transitivity makes $a_1,a_3$ comparable in the
same direction.  The resulting structure $Q$ is a path of
four copies of $D$, with three successive singleton gluing
vertices.  Its two endpoint vertices are $a_1,a_3$.

We verify the essential extra assertion
$\{a_1,a_3\}\le_d Q$.  Its predimension is $4$.
For $U\supseteq\{a_1,a_3\}$ write $U_i=U\cap D_i$.
Predimensions add over the four blocks, subtracting $2$ for
each gluing vertex that is present.  If all three gluing
vertices are present, every block contains its two distinguished
strong singleton points.  Thus $\delta(U_i)\ge3$ and
$\delta(U)\ge12-6=6$.

If a gluing vertex is absent, split the path at all absent
gluing vertices, and discard empty blocks.  The two endpoint
components are distinct.  A component consisting of one block
and containing an endpoint has predimension at least $2$,
with equality only for that bare endpoint, by strongness of
its singleton.  An endpoint component consisting of $r\ge2$
blocks has lower bound
\[
 3+2+3(r-2)-2(r-1)=r+1\ge3:
\]
its first block contains the endpoint and a gluing vertex,
its last contains a gluing vertex, and all intervening blocks
contain two.  A nonempty component containing neither endpoint
has positive predimension.  For $r=1$ this is positivity in
$D$; for $r\ge2$ the corresponding bound is
$2+2+3(r-2)-2(r-1)=r>0$.
Consequently $\delta(U)>4$ unless $U$ consists of the two
bare endpoints.  This proves the asserted strongness.

All ordered strong pairs are conjugate: they are the same
two-point edgeless ordered structure.  Hence every ordered
strong pair is comparable in the direction just found.
For arbitrary $x<y$ in $\Omega$, take a finite strong hull
of them and freely adjoin an isolated strong point $c$ in the
order gap $(x,y)$.  Both $\{x,c\}$ and $\{c,y\}$ are strong:
this follows by adding the isolated point to the strong
singleton $x$, respectively $y$.  Homogeneity and transitivity
give $x\prec c\prec y$, or the reverse throughout.
Thus the entire relation is the given linear order or its
reverse.  If there was no strict comparison, it is equality.
\end{proof}

\begin{proposition}[Proved: separated positive incidence families]
\label{G:separation}
Let $M:A\to\FreeMon$ have injective transitions and let
$u,v\in M(V)$.  Suppose the original DF D.4 group relator
holds in $M(W)^{\gp}$:
\[
 u_2u_1^{-1}v_3'\,(v_1')^{-1}=1.
 \tag{G.12}
\]
In the smooth inclusion-colimit monoid $S$, let $u_x,v_x$
be the corresponding elements at $x\in\Omega$.  Either both
families are constant, or, for any two distinct points $x,y$,
the elements $u_x,u_y$ are incomparable both as left and as
right factors, as are $v_x,v_y$, and as are $u_x,v_y$.

In particular, neither cross-family lift
\[
 v_3'=u_1t,\quad v_1'=u_2t,
 \qquad\hbox{or}\qquad
 u_1=v_3't,\quad u_2=v_1't
 \tag{G.13}
\]
can extend injectively to $A$ while keeping the families
nonconstant.  The same assertion holds with a common left
factor instead of the displayed common right factor.
\end{proposition}
\begin{proof}
An element $u_x$ is fixed by $G_x$.  By (Op) its stabilizer
is $G_F$ for a finite strong $F\subseteq\{x\}$, using
$\operatorname{Fix}(G_x)=\{x\}$.  Thus the family is either
constant or injective with $\operatorname{supp}(u_x)=\{x\}$.
The same holds for $v$.  If one family is constant, (G.12)
makes two distinct-point values of the other equal.  The
other is therefore constant as well.

Assume both are nonconstant.  Prefix order on the $u$-family
pulls back to a $G$-invariant partial order on $\Omega$ with
linearly ordered principal downsets: two prefixes of one word
are comparable in any free chart containing their factorizations.
If any two distinct values are comparable, Lemma
\ref{G:treeorder} makes this order the raw order or its reverse.
Apply the eleven-vertex configuration (G.10).  Put its finitely
many $u$-values, $v$-values and the required nonidentity prefix
quotients in a later free chart of the original inclusion
diagram.  Such a chart exists by filteredness.  Equation
(G.12) holds in its free group, so word lengths give the same
two identities (G.11).  But all six differences
$U_b-U_a$ have the same strict sign, contradicting their
positive sum $\tau_1+\tau_2+\tau_4+\tau_5=0$.
Thus the $u$-family is a prefix antichain.  Suffix order has
the same tree-order property, so exactly the same argument
makes it a suffix antichain.  Applying the argument to $v$
uses $V_a-V_c$ in (G.11) and gives both conclusions for $v$.

Suppose next that $u_x$ is a prefix of $v_y$ with $x\ne y$.
Choose a finite strong $D$ containing $x,y$ and the factor
witness.  Two free copies of $D$ over $y$ give two distinct
$u$-values which are prefixes of the same $v_y$.  They are
comparable, contrary to what was just proved.  If $v_y$ is a
prefix of $u_x$, amalgamate over $x$ and interchange the two
families.  Reversing all words gives the identical argument
for suffixes.  This proves the separation assertion and
excludes (G.13), including $t=1$ in the nonconstant case.
For common left factors, use the suffix conclusion.
\end{proof}

This separation argument alone does not exclude every finite
positive realization of (G.12). Theorem~\ref{G:noincidencelift}
below will exclude its nonconstant global extension.
For example, in the free monoid on $a,b,c$ the four words
\[
 u_1=ac,\quad u_2=bc,\quad
 v_3'=aacca,\quad v_1'=bacca
 \tag{G.14}
\]
are pairwise incomparable both as prefixes and as suffixes,
and satisfy (G.12).  They have the rectangle factorization
$u_1=pq$, $u_2=rq$, $v_3'=ps$, $v_1'=rs$ with
$p=a$, $r=b$, $q=c$, $s=acca$.
Their exponent vectors span $\Z^3$: the determinant of the
columns for $u_1,u_2,v_3'$ is $-1$.
Adding independent letters for $u_3$ and $v_2'$ gives an
injective free-monoid diagram on the full subcategory with
objects $V,W$, whose abelianization is exactly the incidence
cokernel there.  The exact script
\texttt{code/G\_rectangle\_local.py} checks the free-group
reduction, both code conditions, and the unimodular minor.
These finite rectangle data remain valid. The stronger
arbitrary-orbit tree-order theorem below proves that no
nonconstant realization of the same group relator extends
to all strong structures; see Theorem~\ref{G:noincidencelift}.

\begin{proposition}[Proved: a local rectangle realizes the full free-group diagram]
\label{G:localrectangle}
On the full subcategory of $A$ with objects $V,W$ there is an
injective free-monoid diagram whose group completion is the
restriction of DF D.4 and whose abelianization is the incidence
cokernel of R3.  All six point-generator images are a bifix
code, and the incidence difference has both a positive and a
negative coordinate in the free abelian basis.
\end{proposition}
\begin{proof}
Take $M(V)=\FreeMon(u,v)$ and
$M(W)=\FreeMon(a,b,c,d,e)$, with point maps
\[
\begin{array}{c|cc}
 &u&v\\ \hline
1&cbabbaab&cabbaab\\
2&caabbaab&e\\
3&d&cbbbaab
\end{array}
\tag{G.15}
\]
The six words are pairwise incomparable as prefixes and
suffixes, directly from their first and last differing letters;
in particular all three point maps are injective.
Put $s=bbaab$, $p=cb$, $r=ca$, $q=as$.
Then the four nontrivial words are the rectangle
$pq,rq,ps,rs$, so (G.12) holds.
The difference $[u_2]-[u_1]=[a]-[b]$ is mixed, and the
determinant of the exponent columns for $u_1,u_2,v_3'$ on
$a,b,c$ is $-1$.

It remains to check the stronger group assertion, since a
unimodular exponent matrix alone would not imply it.  Set
$x=u_1$, $y=u_2$, $z=v_3'$ and $w=yx^{-1}z=v_1'$.
Within the subgroup they generate,
\[
 A=yw^{-1}=cac^{-1},\qquad
 t=wz^{-1}=cab^{-1}c^{-1},\qquad
 B=t^{-1}A=cbc^{-1}.
\]
As $z=cb^3a^2b$, we recover
\[
 c=(B^3A^2B)^{-1}z,\qquad
 a=c^{-1}Ac,\qquad b=c^{-1}Bc.
\]
Thus $x,y,z$ generate the free group on $a,b,c$; three
generators of a rank-three free group form a basis.
Together with $d,e$ they give precisely the five free
generators of the single-edge DF presentation, with
$v_1'=u_2u_1^{-1}v_3'$ as its only defining relation.
Abelianization gives exactly the R3 cokernel.

The script \texttt{code/G\_rectangle\_full\_group.py}
also expands the displayed inverse words and checks all six
compositions in both directions by free reduction.  It checks
the code condition and integral determinant independently;
Stallings folding is a secondary check, not the sole evidence
for the group isomorphism.  This is a construction only on
$V,W$.  No extension to all of $A$ is asserted.
\end{proof}

\subsection{A general bounded-dimension theorem for total indices}

\begin{proposition}[Proved: module injectivization preserves derived colimits]
\label{G:moduleinjectivization}
Let $J$ be any total category and let $R$ be a ring.  For a
diagram $N:J\to R\text{-}\mathrm{Mod}$ put
\[
 K(a)=\{x\in N(a):N(f)x=0\text{ for some }f:a\to b\},
 \qquad N^{\mathrm{inj}}(a)=N(a)/K(a).
\]
Then $K$ is a subfunctor, all transition maps of
$N^{\mathrm{inj}}$ are injective, and the quotient induces
\[
 L_i\colim_JN\ \cong\ L_i\colim_JN^{\mathrm{inj}}
 \qquad(i\ge0).
 \tag{G.16}
\]
No finite-hom-set or finite-generation assumption is needed.
\end{proposition}
\begin{proof}
Totality implies Condition I.  Two arrows out of the same
object have a common commutative extension.  Thus finitely
many eventual-kernel witnesses can be synchronized by one
arrow.  This proves closure of $K(a)$ under addition and
scalar multiplication.  The same square construction proves
functoriality.  If the image of $x$ is eventually killed after
an arrow $a\to b$, then $x$ is killed by the composite; hence
the quotient transitions are injective.

Let $F$ be the free $R$-module functor on sets and $U$ the
underlying-set functor.  Use the standard augmented simplicial
free-module resolution of $N$, whose degree $n$ is
$(FU)^{n+1}N$.  It is objectwise a resolution: the underlying
augmented simplicial set has its usual extra degeneracy.
Every term has the form $R[L]$ for a set action $L$ on $J$.
Its colimit bar complex is the chain complex of
$\El_J(L)$, with coefficients $R$.  By totality all its
components are weakly contractible, so
$L_i\colim_JR[L]=0$ for $i>0$.
Consequently the homology of the simplicial colimit of this
resolution computes all the derived colimits of $N$.

The map to the analogous resolution of $N^{\mathrm{inj}}$
induces an isomorphism after colimit in every degree.
Surjectivity follows pointwise.  To prove injectivity, lift a
finite formal expression and a proposed equality between its
images.  Each iterated free-module term consists of finite
formal expressions, nested a finite number of times.  At
their leaves, equality after quotient says that finitely many
differences of elements of $N(a)$ are eventually killed.
Synchronize those finitely many witnesses.  At that common
extension the leaf equalities already hold, hence so does
the entire equality of formal expressions.  If a colimit
equality uses a zigzag, lift its finitely many expressions
and apply this argument at each step; the arrows of the
zigzag themselves also lift by naturality.  Thus the map
on colimits is injective.  The resulting simplicial colimit
complexes are isomorphic, proving (G.16).
\end{proof}

\begin{proposition}[Proved: bounded dimensions are acyclic on every total index]
\label{G:boundedcoefficients}
Let $J$ be total, let $k$ be any field, and let
$N:J\to k\text{-}\mathrm{Mod}$ satisfy
$\dim_kN(a)\le d$ for one finite integer $d$ and every $a$.
Then $L_i\colim_JN=0$ for all $i>0$.
Moreover the kernel of $N(a)\to\colim_JN$ is exactly $K(a)$
from Proposition~\ref{G:moduleinjectivization}.

If instead every $N(a)$ is a finitely generated free abelian
group and their ranks have one finite uniform bound, the same
acyclicity conclusion holds, and $\colim_JN$ is torsion-free.
In particular the abelianization of every uniformly bounded
finite-rank free-monoid diagram is acyclic over every total
index, even without the Fra\"iss\'e hypotheses of this part.
\end{proposition}
\begin{proof}
Replace $N$ by its injectivization, using (G.16).  Its dimensions
remain bounded.  Let $d_0$ be the largest attained dimension,
and let $D\subseteq J$ be the full subcategory where it is
attained.  Since all transitions are injective, $D$ is a
cosieve and its transition maps are isomorphisms.

A connected category satisfying I has common successors:
successively complete the spans along a zigzag connecting
two objects.  Thus every object of $J$ maps to an object of
$D$.  For $a\in J$, the category $a/D$ is a nonempty full
cosieve in $a/J$, and Condition I makes it connected.
It is the category of elements of the subaction of
$J(a,-)$ consisting of arrows whose targets lie in $D$.
Totality therefore makes $a/D$ weakly contractible.
So $D\hookrightarrow J$ is homotopy final.
Likewise $D$ itself is the category of elements of a
subterminal action, is connected by common successors, and
has weakly contractible nerve.

The restricted vector-space diagram has only isomorphisms,
so it factors through the fundamental groupoid of $D$.
That groupoid is equivalent to the terminal groupoid, because
$BD$ is connected and simply connected.  The restriction is
therefore isomorphic to a constant diagram.  Its positive
derived colimits vanish.  Homotopy finality gives the same
over $J$.  It also shows that every map from an injectivized
value to its colimit is injective: map the object into $D$,
where its transition is injective and the maps to the
constant colimit are isomorphisms.  This proves the kernel
assertion for the original $N$.

For the integral assertion, use the colimit bar complex
$C_\bullet$ of $N$.  Each term is free abelian.  The field
case applies to $N\otimes\Q$ and to $N\otimes\F_p$ for
every prime $p$.  Therefore $H_i(C_\bullet)\otimes\Q=0$
for $i>0$.  The universal-coefficient sequence for
$C_\bullet\otimes\F_p$, in degrees $i$ and $i+1$, shows
both $H_i(C_\bullet)/pH_i(C_\bullet)=0$ and
$H_i(C_\bullet)[p]=0$ for every $i>0$.
The first rational conclusion says these groups are torsion;
the absence of prime-order elements then makes them zero.
The degree-one universal-coefficient sequence also gives
$H_0(C_\bullet)[p]=0$ for every $p$, so the ordinary colimit
is torsion-free.  This proves all assertions.
\end{proof}

\begin{proposition}[Proved: injective colimit cocones and all rank-one diagrams]
\label{G:generalcomparison}
Let $J$ be any total category and $M:J\to\Mon$.
If every canonical map
$M(a)\to C=\colim_J^{\Mon}M$ is injective, then the specified
canonical comparison $q_M$ is a weak equivalence.
Moreover every diagram whose values have free rank at most
one passes over every total $J$, including diagrams with
erasing transitions.
\end{proposition}
\begin{proof}
Under the first hypothesis, identify $M(a)$ with its image
submonoid in $C$, and let $P$ be the poset of these image
submonoids, ordered by inclusion.  Common successors in $J$
make $P$ filtered.  Its union is $C$: a finite product of
elements from several values can be moved into one common
successor.  The diagram factors through the inclusion diagram
on $P$.  For $P_0\in P$, the comma category of the functor
$J\to P$ under $P_0$ is the full cosieve of objects $a$ for
which $P_0\subseteq M(a)$.  It is nonempty and connected,
again by common successors.  As the category of elements of
a subterminal action on $J$, it is weakly contractible by
totality.  Hence $J\to P$ is homotopy final.  The comparison
is the composite
\[
 \operatorname{hocolim}_{a\in J}BM(a)
 \xrightarrow{\sim}
 \operatorname{hocolim}_{P_0\in P}BP_0
 \xrightarrow{\sim} BC.
\]
The second arrow is the filtered-colimit comparison; the
nerve of a monoid commutes with filtered colimits, degree by
degree, since finite products do.  This proves the first
assertion for the specified canonical map.

For the second assertion first quotient $M(a)$ by eventual
equality: two elements are equivalent when one arrow out of
$a$ makes them equal.  Condition I synchronizes witnesses,
so this is a monoid congruence, gives a diagram $M^{\mathrm{inj}}$
with injective transitions, and leaves the ordinary monoid
colimit unchanged.  This quotient also leaves the homotopy
colimit of classifying spaces unchanged.  Indeed for each
$n\ge0$ the induced map of set actions
$M^n\to(M^{\mathrm{inj}})^n$ gives a bijection of ordinary
set colimits: it is pointwise surjective, and equality of
finitely many coordinates can be witnessed simultaneously
after one arrow.  Both homotopy colimits of these set actions
have weakly contractible components by totality, so that map
is a weak equivalence in each degree of the monoid nerve.
Realizing the simplicial homotopy-colimit comparison proves
the assertion.  The square of canonical comparisons commutes,
so passage transfers back to $M$.

A rank-one monoid is $(\N,+)$.  Its homomorphisms to another
rank-one monoid are multiplication by nonnegative integers.
Such a map is either injective or zero.  Thus after the
quotient every value remains either $1$ or $\N$.
If all are $1$, the result is $BJ\simeq *$.
Otherwise the nontrivial values form a nonempty cosieve $D$.
It is homotopy final: $a/D$ is a nonempty connected cosieve
of $a/J$, and is an action component, exactly as in
Proposition~\ref{G:boundedcoefficients}.
The category $D$ is itself total, because an action on a
cosieve extends by the empty set outside that cosieve; it is
connected by common successors.

On $D$, the transition multipliers lie in $\N_{>0}$.
Viewed in the group $\Q_{>0}^{\times}$ they form a local
system on the weakly contractible nerve $BD$.
Choose positive rational numbers $\lambda_a$ satisfying
$\lambda_b n_f=\lambda_a$ for every $f:a\to b$.
Then $r\mapsto\lambda_a r$ identifies the diagram with
inclusions of the submonoids $\lambda_a\N$ of $\Q_{\ge0}$.
Their images are directed by common successors.  Their union
is their monoid colimit: equal rational representatives
become equal in a common successor, so no additional
identifications occur.  The canonical maps are injective,
and the first assertion applies.  Homotopy finality and the
preceding quotient comparison finish the rank-one case.
\end{proof}

\subsection{Gap diagrams on all lexicographic embedding actions}

\begin{theorem}[Proved: lexicographic embedding-gap diagrams pass]
\label{G:embeddinggaps}
Fix a nonempty $F_0\in\Cdelta$ and any fixed listing of its
vertices.  Order $L(X)=\Cdelta(F_0,X)$ lexicographically by
that listing; the increasing listing is a special case.
Let $M_{F_0}(X)$ be the free
monoid on the consecutive gaps of
$\{-\infty\}\cup L(X)\cup\{+\infty\}$, with a strong
embedding acting by subdivision of gaps.  The set $L(X)$
may be empty.
Then the specified comparison $q_{M_{F_0}}$ is a weak
equivalence and $T_i((M_{F_0})_{\ab})=0$ for every $i>0$.
Its target is a $K(F_2,1)$.

More precisely, list the distinct strong closures in $F_0$
of the initial segments of the chosen listing as
\[
 \varnothing=B_0\subsetneq B_1\subsetneq\cdots
 \subsetneq B_d=F_0.
\]
The ordinary monoid colimit has presentation
\[
 P_d=\langle l,r,e_0,\ldots,e_{d-1}\mid
 e_i e_j=e_{\min(i,j)},\quad le_i=l,\quad e_i r=r
 \ (0\le i,j<d)\rangle.
 \tag{G.17}
\]
In particular the strong-hyperedge gap diagram, $F_0=W$,
has two distinct nonidentity idempotents $e_0,e_1$ in its
ordinary colimit and still passes.
\end{theorem}
\begin{proof}
Write $\Omega=\operatorname{Emb}_{\le}(F_0,M_0)$ with this
lexicographic order.  The inclusion-colimit monoid is the
interval monoid on $\Omega$ with its two formal endpoints:
every finite collection of embeddings lies in some finite
strong substructure.  Its orbit monoid is the ordinary
colimit, as in Proposition~\ref{G:saturation}.
There is one orbit of $\Omega$.

Agreement on $B_k$ defines a convex equivalence relation on
$\Omega$.  Indeed $B_k$ is the strong closure of an initial
segment, and homogeneity and uniqueness of strong closure
show that agreement on that segment is equivalent to agreement
on its closure.  Explicitly, an automorphism fixing the segment
preserves its finite strong closure setwise, hence pointwise
because of the linear order.
Equality of an initial segment is a convex
condition in lexicographic order.  For distinct $x<y$ in
$\Omega$, call the largest $k$ for which their embeddings
agree on $B_k$ the depth of the pair.

We show that the orbit class of an interval $[x,y]$ depends
only on this depth.  Fix depth $k$ and identify the common
copy of $B_k$.  The first vertex $j$ in the chosen listing
outside $B_k$
is exactly the first coordinate where $x,y$ differ: agreeing
there would force agreement on $B_{k+1}$, its next closure.
In particular $x_j<y_j$.  Choose a finite strong $D$ containing
both copies of $F_0$, and freely amalgamate two new copies
$z,t$ of $F_0$ with $D$ over their common $B_k$.

Choose the raw order in each gap of $B_k$ as follows.
In the gap containing the pivot coordinates $x_j,y_j$,
put all new vertices of $z$ in a short interval just after
$x_j$, before $y_j$ and before the next raw new vertex of
$x$, if one exists.  Put all new vertices of $t$ just before
$y_j$, after the preceding raw new vertex of $y$, if one
exists, and after the entire $z$ cluster.  There is room for
both choices because $x_j<y_j$; use the boundaries of the
gap when a neighboring raw vertex is absent.
In every other gap containing new vertices, put the $z$
cluster, then the $t$ cluster, before the new vertices of
both old copies.  Within each cluster keep the prescribed
raw order of $F_0$.
These choices are permitted ordered free amalgams and give
$x<z<t<y$ in the chosen lexicographic order.

All three relevant two-copy unions are strong free amalgams
over $B_k$.  The ordered pairs $(x,z)$ and $(t,y)$ have
fixed isomorphism types: in the pivot gap their types depend
only on the raw rank of the pivot vertex, and in the other
gaps the new copy precedes the old.  Denote their interval
classes by $a_k,b_k$.  The pair $(z,t)$ has the fixed type
in which the first copy's new vertices precede the second
copy's in every raw gap of $B_k$; denote its class by $e_k$.
Thus every interval of depth $k$ has the same class
$a_ke_kb_k$.
Apply the construction when the original pair already has
that separated type.  It gives $e_k=a_ke_kb_k$.
All intervals of depth $k$ therefore have class $e_k$.

For $x<y<z$, convexity of the nested equivalences gives
\[
 \operatorname{depth}(x,z)=
 \min\{\operatorname{depth}(x,y),
          \operatorname{depth}(y,z)\}.
\]
Every pair of depths is realized by taking free copies over
the corresponding nested strong bases.  Subdivision yields
the multiplication relations in (G.17).  The two ray classes
are $l,r$, and their subdivisions give $le_k=l$ and $e_kr=r$.
Conversely assign to an interval its $e_k$, to each ray its
specified $l$ or $r$, and to the whole interval $lr$.
The displayed relations check every subdivision, including
the empty finite set.  This assignment is invariant under
$G$ and is inverse to the presentation map, proving the
full monoid presentation.  Its idempotents $e_k$ are distinct
and nonidentity: send $l,r$ to $e_0$ in the finite chain
semilattice with these distinct elements and a separately
adjoined identity.

For the comparison, let $\mathcal E_d$ be the category of
finite linearly ordered sets with $d-1$ nested convex
equivalence relations, including the empty object, with
increasing induced embeddings.  Its Fra\"iss\'e limit is
$\Q^d$ in lexicographic order, with equivalences given by
agreement of successive initial coordinate strings.
Universality and homogeneity follow recursively by embedding
the finitely many blocks in $\Q$ and extending finite order
isomorphisms independently inside those blocks.
Its automorphism group is an iterated semidirect product
\[
 H_1=\Aut(\Q,<),\qquad
 H_d=\Aut(\Q,<)\ltimes H_{d-1}^{\Q}.
\]
Finite-order Ramsey gives extreme amenability of $H_1$;
products and extensions preserve extreme amenability by the
compact fixed-set argument of KPT Lemma~6.7.  Induction and
KPT Theorem~4.8 give the Ramsey property of $\mathcal E_d$
\cite{ABKPT2005}.  For $d=2$ this is their Corollary~6.8.
The empty object adds no Ramsey constraint.  Hence
$\mathcal E_d$ is exact by R5.

The functor $X\mapsto L(X)$, with the agreement relations
on the $B_k$, takes values in $\mathcal E_d$ and pulls its
gap diagram back to $M_{F_0}$.  In the template, its two-point
objects identify intervals by depth directly, and its
three-point objects impose precisely the minimum-depth
subdivision rule.  The resulting finite-interval presentation
is (G.17) as well;
the induced map of colimits is the identity on all displayed
generators.  Natural prefix bases identify the group-completed
diagram with the free-group diagram on $L(X)\sqcup\{*\}$.
The graph argument of Theorem~\ref{G:gap} applies: totality
contracts the element-category component of the transitive
embedding action, and the constant generator supplies the
second loop.  Both homotopy colimits are $K(F_2,1)$, and
the template map identifies their two free generators.
It is a weak equivalence.  The template's specified comparison
passes by Proposition~\ref{G:exact}; the commuting comparison
square now proves passage for $M_{F_0}$.
Finally the same prefix calculation gives
$(M_{F_0})_{\ab}\cong\Z[L]\oplus\Z$, whose positive
derived colimits vanish by totality.
\end{proof}

\subsection{Intrinsic locally free tests}

\begin{proposition}[Proved: bounded rank forbids absorption in the colimit]
\label{G:boundedtarget}
Let $J$ be any total category, and let $M:J\to\FreeMon$ have
free ranks bounded by one finite integer $d$. Put $C=\colim_JM$
in $\Mon$, and let $C_{\ab}$ denote its universal abelian group.
Then $C_{\ab}$ is torsion-free of rank at most $d$, and the map
$C\to C_{\ab}$ has inverse image of $0$ equal to $\{1\}$.
Moreover, the image $P\subseteq C_{\ab}$ satisfies
$P\cap(-P)=\{0\}$. In particular $C$ has no nonidentity
idempotent or unit, and
\[
 cx=c\ \text{or}\ xc=c\quad\Longrightarrow\quad x=1.
\]
No Fra\"iss\'e hypothesis is needed.
\end{proposition}
\begin{proof}
The universal abelian group functor preserves colimits, so
$C_{\ab}=\colim_JM_{\ab}$. Proposition~\ref{G:boundedcoefficients}
gives torsion-freeness. Any finitely many elements of this colimit
can be represented at one object, since $J$ has common successors.
Their rational span therefore has dimension at most $d$.

Every $c\in C$ is also represented by one word $w\in M(a)$:
a finite product of stage elements can be moved to a common
successor. If $[c]=0$, the kernel assertion of
Proposition~\ref{G:boundedcoefficients}, applied over $\Q$, gives
an arrow $f:a\to b$ for which the exponent-count vector of
$M(f)(w)$ is zero. All its coordinates are nonnegative integers,
so $M(f)(w)=1$ and $c=1$. If $[c]+[x]=0$, choose simultaneous
word representatives at one object and apply the same argument
to their product. Its eventual image is empty, so each factor's
image is empty. Thus $c=x=1$, which proves pointedness of $P$
and the assertion about units. An idempotent or a one-sided
stabilizer has $[x]=0$ in $C_{\ab}$, proving the remaining claims.
This conclusion alone does not assert comparison passage in
bounded rank over an arbitrary total category.
\end{proof}

\begin{lemma}[Proved: functorial free hulls and locally free tests]
\label{G:hullfunctorial}
Call a monoid locally free if every finite subset lies in a
finite-rank free submonoid. If $S,T$ are locally free, $X\subset S$
is finite, and $f:S\to T$ is any monoid homomorphism, then
\[
 f(H_S(X))\subseteq H_T(fX),
\]
where $H$ denotes the intrinsic least free hull. Erasing and
noninjective homomorphisms are allowed.

Let $J$ be a small category with finite hom-sets. Every diagram
$N:J\to\Mon$ with locally free values is a filtered union of
subdiagrams with finite-rank free values. Consequently every
free-monoid comparison over $J$ is a weak equivalence if and
only if the same holds for all locally free monoid diagrams.
\end{lemma}
\begin{proof}
We give the closure argument, including its ambient-monoid issue.
A submonoid $H\subseteq S$ is stable when
\[
 p,q,pw,wq\in H\quad\Longrightarrow\quad w\in H
 \qquad(p,q,w\in S).
 \tag{G.18}
\]
In a free ambient monoid, stability is equivalent to freeness
\cite[Theorem~1, p.~84]{GCarpi2014}. The inductive free-hull
description by the unit, generators, multiplication, and
(G.18) is also given in the opened primary formalization
\cite[\texttt{Submonoids}: \texttt{free\_hull.intros},
\texttt{free\_basis\_code}, \texttt{free\_hull\_min}]{GHolubRaskaStarosta2021}.

Every free submonoid $H$ of a locally free $S$ is stable in $S$.
Indeed apply Levi's lemma inside $H$ to
$(pw)q=p(wq)$. One comparison gives $pw=pt$ with $t\in H$,
so cancellation in $S$ gives $w=t$; the reverse comparison
gives $p=(pw)t$, hence $wt=1$, so conicality gives $w=1$.
These cancellation and conicality properties hold in $S$ because
each finite equality lies in a free chart.

Let $K_S(X)$ be the least submonoid containing $X$ and closed
under (G.18). It is the union of the successive finitary closure
steps. Choose a finite free chart containing $X$. The intrinsic
least free hull $H_S(X)$ exists by the intersection argument of
Proposition~\ref{G:saturation}; that argument uses only locally
free charts. Since this hull is stable in $S$,
$K_S(X)\subseteq H_S(X)$. The submonoid $K_S(X)$ is stable
inside the free monoid $H_S(X)$, hence is free itself. Minimality
of the hull forces equality. A homomorphism carries every instance
of the finitary inference (G.18) to an instance of the same
inference. Induction on closure steps therefore proves the
displayed functoriality. This argument requires no lifting of
factorizations through $f$.

Now choose a finite set $D$ of tagged elements $(a,x)$ with
$x\in N(a)$. At $b\in J$, put
\[
 X_D(b)=\{N(u)x:(a,x)\in D,\ u\in J(a,b)\},
 \qquad N_D(b)=H_{N(b)}(X_D(b)).
\]
The set $X_D(b)$ is finite. Functoriality of the hull makes $N_D$
a subfunctor of $N$, with finite-rank free values. Increasing
$D$ gives inclusions of subfunctors, and their filtered union
is $N$. The construction includes $D=\varnothing$, whose hull
is the trivial monoid.

Filtered colimits commute with finite products of sets, hence
with monoid nerves, and with the bar construction for a fixed
small indexing category. They also preserve weak equivalences
of simplicial sets. The comparison for $N$ is therefore the
filtered colimit of the comparisons for $N_D$. This proves the
claimed equivalence of tests; the converse is immediate because
free monoids are locally free.
\end{proof}

\begin{theorem}[Proved: an intrinsic group criterion for the free-monoid test]
\label{G:groupcriterion}
Assume setting (S), and assume its strong embedding category $A$
is total. For a smooth action of $G=\Aut(M_0)$ by automorphisms
on a locally free monoid $S$, let $(NS)/G$ be the degreewise
orbit simplicial set, with $n$-simplices $S^n/G$. Let $S_G$ be
the monoid quotient by the congruence generated by $s=gs$.
Then $A\in\mathcal F$ if and only if, for every such $S$, the
canonical map
\[
 \theta_S: |(NS)/G|\longrightarrow B(S_G)
 \tag{G.19}
\]
is a weak equivalence.

For each $i>0$, the assertion that
$L_i\colim_A M_{\ab}=0$ for every $M:A\to\FreeMon$ is
equivalent to
\[
 H_{i+1}(|(NS)/G|;\Z)=0
 \quad\text{for every smooth locally free }G\text{-monoid }S.
 \tag{G.20}
\]
Consequently both the class-$\mathcal F$ property and these
abelianization-acyclicity properties are invariant under changing
the total Fra\"iss\'e presentation of the same topological group.
Together with Corollary~\ref{A:totality-group-invariant}, membership
in $\mathcal F$ depends only on the topological group throughout
setting (S), without the extra hypotheses (Rig), (SA), (FA),
(FA$\varnothing$), (St), (Op), or (WEI).
\end{theorem}
\begin{proof}
We work with the equivalent category $A_{M_0}$ of finite nonempty
strong subsets of the limit. Let $\mathcal I$ be its directed
poset of inclusions. For any monoid diagram $M:A_{M_0}\to\Mon$,
put $S=\colim_{\mathcal I}M$. An automorphism $g\in G$ sends
the class represented by $m\in M(B)$ to the class represented
by $M(g|_B)m\in M(gB)$. This is well defined and smooth.
Every embedding between finite strong sets extends to such a $g$.
It follows, by presenting a colimit as generators and transition
relations, that
\[
 \colim_{A_{M_0}}\Und(M)^n=S^n/G\quad(n\ge0),
 \qquad \colim_{A_{M_0}}^{\Mon}M=S_G.
 \tag{G.21}
\]
For the first identity, filtered colimits over $\mathcal I$
commute with the finite power; inclusion relations form the
filtered colimit, and the remaining embedding relations are
exactly the diagonal $G$-orbit relations. For the second,
multiplication is already defined in the filtered colimit, and
the remaining relations generate the stated monoid congruence.
In particular the orbit set and monoid coinvariants are not
being identified.

Totality gives a canonical weak equivalence from the homotopy
colimit of each discrete set diagram to its set colimit: its
source is the nerve of the category of elements, and every
component is weakly contractible. Apply this to $\Und(M)^n$
for every $n$. Realizing the resulting degreewise weak
equivalences in the bisimplicial bar construction gives a
natural weak equivalence
\[
 \hocolim_{A_{M_0}}BM\simeq |(NS)/G|.
 \tag{G.22}
\]
Under (G.21)--(G.22), the specified comparison is precisely
$\theta_S$. This establishes the needed saturation comparison
here directly from totality, with no functorial-amalgamation
assumption.

If $M$ has free values, $S$ is locally free. To check this for
possibly noninjective transition maps, represent a finite set
at one stage and minimize the rank of a finite free chart
containing its images over all later stages. A later failure
of injectivity on that chart would lower the rank by the
defect theorem, contradicting minimality. Thus the chosen
chart embeds in $S$, as in Proposition~\ref{G:saturation}.
Consequently (G.19) for all $S$ implies the free-monoid test
over $A$.

Conversely, start with an arbitrary smooth locally free
$G$-monoid $S$. Set $T_S(B)=S^{G_B}$. For a strong embedding
$u:B\to C$, choose an extension $g\in G$ and send $s$ to $gs$.
This is independent of the extension, and its image is fixed
by $G_C$, so it defines a monoid diagram. A finite set fixed
by $G_B$ has its intrinsic free hull fixed pointwise: the
group preserves the hull and permutes its atoms, and every
atom occurs in one of the fixed words. Hence each $T_S(B)$
is locally free. Smoothness gives
$\colim_{\mathcal I}T_S=S$: an open element stabilizer
contains the pointwise stabilizer of a finite strong set.
The hom-sets of $A$ are finite, since its structures are finite.
If $A\in\mathcal F$, Lemma~\ref{G:hullfunctorial} therefore
makes the comparison for $T_S$ a weak equivalence. Applying
(G.21)--(G.22) to this diagram proves (G.19).

For (G.20), a locally free monoid $L$ is a filtered union of
finite free submonoids. Therefore $H_q(BL;\Z)=0$ for $q>1$,
and $H_1(BL;\Z)=L_{\ab}$. Since $BA$ is contractible, the
two-row homology spectral sequence gives natural isomorphisms
\[
 H_{i+1}(\hocolim_A BN;\Z)
     \cong L_i\colim_A N_{\ab}\qquad(i>0)
 \tag{G.23}
\]
for every locally free diagram $N$. If the asserted vanishing
holds for all free diagrams, write $T_S$ as the filtered union
in Lemma~\ref{G:hullfunctorial}. Abelianization and the bar
complex commute with this filtered colimit, so the vanishing
holds for $T_S$; (G.22)--(G.23) give (G.20). The reverse
direction applies the same identities to an arbitrary free
diagram and its locally free filtered limit $S$.

The tests (G.19) and (G.20) use only the topological group:
a topological group isomorphism transports smooth actions,
orbit simplicial sets, and monoid coinvariants. This proves
presentation independence among total presentations. Finally
$\mathcal F\subseteq\mathrm{total}$, and totality itself is
group invariant by Corollary~\ref{A:totality-group-invariant}.
Thus if either presentation belongs to $\mathcal F$, both are
total and the preceding criterion applies; if neither is
total, neither belongs to $\mathcal F$.
\end{proof}

The intrinsic criterion narrows the remaining Part G questions;
it does not decide whether (G.19) always holds or whether
(G.20) can fail for the particular group of $\Cdelta$.

\subsection{Orders on arbitrary support orbits and the incidence relator}

\begin{theorem}[Proved: invariant tree orders are finite signed lexicographic orders]
\label{G:alltreeorders}
Let $F\in\Cdelta$ be finite and strong, and put
$\Omega_F=\Cdelta(F,M_0)$. Let $\preceq$ be a $G$-invariant
partial order on $\Omega_F$ whose principal downsets are chains.
Then there are strong subsets and coordinates
\[
 E=B_0\subsetneq B_1\subsetneq\cdots\subsetneq B_d=F,
 \qquad a_i\in F\setminus B_{i-1},\quad
 B_i=\operatorname{cl}_F(B_{i-1}\cup\{a_i\}),
\]
and signs $\epsilon_i\in\{+1,-1\}$ with the following description.
Two distinct embeddings are comparable exactly when they agree
on $E$. Within each such fiber their order is lexicographic:
at the first $i$ for which their restrictions to $B_i$ differ,
compare their values at $a_i$, in the raw order if
$\epsilon_i=+1$ and in its reverse if $\epsilon_i=-1$.
The equality order corresponds to $E=F$ and $d=0$.

The same assertion holds relative to a fixed embedding of any
strong $B\subseteq F$, with $G_B$ in place of $G$, and then
$B\subseteq E$. In particular every comparability component
of such an order is a chain.
\end{theorem}
\begin{proof}
We use strict comparison $x\prec y$ for $x\preceq y$, $x\ne y$.
The proof below also works with fixed strong parameters $B$;
all amalgams and automorphisms then fix those parameters.

\emph{A comparison gives a strong free pair over its equalizer.}
Suppose $x\prec y$, and choose a finite strong $D$ containing
$xF\cup yF$. Let $E\subseteq F$ be the coordinate equalizer,
the set of $a$ with $x(a)=y(a)$. The overlap induces an
increasing partial isomorphism $\sigma$ of $F$, with strong
domain and range. Its fixed points are $E$. Let
$C=\operatorname{cl}_F(E)$. The strong domain and range both
contain $C$, and $\sigma(C)$ is strong and contains $E$.
Minimality gives $C\subseteq\sigma(C)$; equal finite cardinalities
give equality. The order on $C$ makes $\sigma|_C$ the identity,
so $C=E$. Thus $E$ is strong.

Concatenate copies of $D$, freely amalgamating the second
$F$-copy of each block with the first of the next by their
coordinate identification. These are ordered strong amalgams,
and consecutive endpoint embeddings remain strictly comparable.
An endpoint coordinate can persist through $m$ blocks only if
the corresponding $m$-fold partial iterate of $\sigma$ is defined.
A nonfixed increasing partial-bijection orbit in a finite chain
is strictly monotone, so it dies within $|F|$ steps. Take at
least $|F|$ blocks. Their two endpoint $F$-copies therefore
meet precisely in the common image of $E$. Call this longer
block $D'$.

Write $\delta_E(X)=\delta(X)-\delta(E)$ for $E\subseteq X$
and put $d_0=\delta(F)-\delta(E)>0$. Concatenate $r>d_0$
copies of $D'$ in the same fashion. Denote their successive
endpoint copies by $F_0,\ldots,F_r$. Adjacent endpoint copies
meet only in $E$; nonadjacent blocks also meet only in $E$.
For any induced $X$ containing $F_0\cup F_r$, additivity over
the free path gives
\[
 \delta_E(X)=\sum_{i=1}^{r}\delta_E(X\cap D'_i)
       -\sum_{i=1}^{r-1}\delta_E(X\cap F_i).
 \tag{G.24}
\]
If $X$ meets every internal $F_i$ outside $E$, the first block
strictly enlarges its strong left endpoint, and each later
block adds a point outside its incoming strong $F_{i-1}$.
The latter statement uses $F_{i-1}\cap F_i=E$.
Submodularity and strict strongness give an increment at
least one even when $X$ contains only part of a glue:
\[
 \delta(X\cap D'_i)-\delta(X\cap F_{i-1})\ge1.
\]
Hence $\delta_E(X)\ge d_0+r>2d_0$.
If an internal glue is empty outside $E$, cut the path at all
such glues. Each endpoint component has relative predimension
at least $d_0$, with equality only for its bare endpoint $F$,
because that endpoint is strong in the whole path. Every
other nonempty component has positive relative predimension,
since $E$ is strong. Thus every strict enlargement of
$F_0\cup F_r$ has relative predimension greater than $2d_0$.
The endpoint union is therefore strong. Since $r\ge2$ it is
the free amalgam of its two $F$-copies over $E$.
Transitivity retains their comparison. We have proved the
claimed free-pair reduction.

\emph{An active base has a single first coordinate.}
Call a strong $E\subsetneq F$ active if some strong free
$F$-pair over $E$, agreeing there, is comparable. Fix an active
$E$ and an oriented witness $x\prec y$. For each coordinate
outside $E$, choose a small raw-order interval around its
value in $x$, with these intervals mutually disjoint and
avoiding all coordinates of $y$ outside $E$. Two new copies
of $F$ in these boxes, freely amalgamated over $E$, each have
the same pair type with $y$ as $x$ does. They both lie below
$y$, so they are comparable. Their mutual pair type is
coordinate-separated: every coordinate of one copy and the
corresponding coordinate of the other lie in one of the
chosen disjoint boxes.

For coordinate-separated free pairs over $E$, the pair type
depends only on which of the two values is smaller at each
coordinate in $F\setminus E$. Thus comparison is a function
of this finite sign pattern. From the local comparison just
obtained, put two lower copies in the same prescribed
coordinate half-boxes below one upper copy. Their mutual
sign pattern can be arbitrary. The chain-downset property
therefore forces comparability for every sign pattern.
Asymmetry gives exactly one orientation for each pair of
opposite patterns.

Here is the elementary finite-coordinate order argument.
Reverse the social comparison if necessary so that unanimous
coordinate increase implies increase. For the finite coordinate
set $I=F\setminus E$, let $\mathcal U\subseteq\mathcal P(I)$
consist of the sets of coordinates at which the first tuple
is smaller when that tuple precedes the second. Exactly one
of $A,A^c$ belongs to $\mathcal U$, and $I\in\mathcal U$.
If $A\subseteq C$ and $A\in\mathcal U$, realize three tuples
whose comparison sign sets for $(x,y),(y,z),(x,z)$ are
$A,I,C$: use the coordinate orders $x<y<z$ on $A$,
$y<x<z$ on $C\setminus A$, and $y<z<x$ outside $C$.
Transitivity gives $C\in\mathcal U$.
If $A,C\in\mathcal U$, realize the corresponding sign sets
$A,C,A\cap C$ using respectively
\[
 x<y<z,\quad z<x<y,\quad y<z<x,\quad z<y<x
\]
on $A\cap C,A\setminus C,C\setminus A,I\setminus(A\cup C)$.
Transitivity gives $A\cap C\in\mathcal U$.
All these profiles are realizable by strong free copies in
the disjoint coordinate boxes. Thus $\mathcal U$ is a proper
ultrafilter on a finite set, hence consists exactly of the
subsets containing one coordinate $a$. Returning to the
original comparison, there is a sign $\epsilon$ such that
every coordinate-separated pair is ordered by its values
at $a$, with that sign.

The same rule holds for every strong free pair over $E$.
Realize its combined raw order by real coordinates, keeping
the coordinates of $E$ fixed. The two tuples lie in the same
open cell of strictly increasing coordinates in the prescribed
$E$-gaps. The straight segment between them stays in that cell.
When their values at $a$ are in the chosen signed order,
subdivide this segment finely enough that successive tuples
are coordinate-separated. Choose the finitely many intermediate
parameters generically to avoid all unintended coordinate
coincidences; each coordinate varies strictly because the
endpoints are disjoint outside $E$. Realize this finite raw
order as a free amalgam of $F$-copies over $E$, over the given
endpoint pair. Each consecutive comparison has the chosen
orientation, so transitivity proves the rule for the endpoints.

Finally let $x,y$ be arbitrary embeddings agreeing on $E$,
with distinct values at the pivot $a$. Take a finite strong
set containing both images, and freely adjoin one $F$-copy
$z$ over $E$, with $z(a)$ strictly between $x(a)$ and $y(a)$.
The pairs $(x,z)$ and $(z,y)$ are strong free pairs over $E$:
free amalgamation preserves the strong subobjects $xF,yF$.
The rule just proved therefore compares $x,z,y$ in order.
Thus every active base gives the global signed pivot rule,
even for pairs with additional overlaps or relations.

\emph{The least active base and the successive pivots.}
Active bases are closed under intersection. For active
$D_1,D_2$, fix an upper $F$-copy $y$. Freely adjoin an
$F$-copy $x_i$ over its restriction to $D_i$, choosing the
signed pivot of $D_i$ below that of $y$. Then $x_i\prec y$.
The two new copies have coordinate equalizer exactly
$D_1\cap D_2$. They are comparable by the chain-downset
property, and the free-pair reduction makes their equalizer
active. There are only finitely many strong subsets of $F$,
so, if any comparison exists, there is a least active base $E$.
The equalizer of every comparable pair is active by the
first step; hence comparable pairs agree on $E$.

Within an $E$-fiber, distinct pivot values have already been
compared. If two distinct embeddings have the same pivot
value, freely adjoin over $E$ an upper copy with its signed
pivot above that common value. Both original embeddings
lie below it, so they too are comparable. Thus each
$E$-fiber is a chain and is exactly a comparability component.

Put $B_1=\operatorname{cl}_F(E\cup\{a\})$.
Agreement on $E$ and $a$ is equivalent to agreement on $B_1$:
the overlap partial isomorphism fixes the former set, and
the finite ordered closure argument used above makes it
fix the latter pointwise. If $B_1\ne F$, restrict to the
embeddings extending one fixed $B_1$-embedding. Their order
is total and $G_{B_1}$-invariant. The argument above, with
$B_1$ as fixed base, supplies the next signed pivot outside
$B_1$. Iterating strictly enlarges the finite strong base,
and terminates at $F$. This gives exactly the asserted
signed lexicographic rule. If there was no strict comparison,
take $E=F$.
\end{proof}

\begin{theorem}[Proved: the full incidence group relator has no nonconstant positive realization]
\label{G:noincidencelift}
Let $S$ be any smooth locally free $G$-monoid for
$G=\Aut(M_0)$ of $\Cdelta$. Let $x\mapsto u_x$ and
$x\mapsto v_x$ be $G$-equivariant maps from the orbit
$\Omega$ of strong singletons to $S$. Suppose that for
every strong ordered hyperedge $a<b<c$,
\[
 u_bu_a^{-1}v_c v_a^{-1}=1\quad\text{in }S^{\gp}.
 \tag{G.25}
\]
Then both families are constant. In particular the DF D.4
free-group diagram has no free-monoid lift over all of
$\Cdelta$ with these nonconstant point generators, even after
allowing arbitrary extra free generators and arbitrary
rectangle factorizations at the hyperedges.
\end{theorem}
\begin{proof}
The separation argument of Proposition~\ref{G:separation}
applies to these data as well. We explain the only extension
of its hypotheses. A locally free $S$ is the filtered union
of its finite free charts, and its group completion is their
filtered group colimit. A finite collection of group
equalities is therefore witnessed in a later finite free
chart, by choosing charts containing all the required words
and then the finitely many equality witnesses. Positive
elements embed in $S^{\gp}$: an equality between positive
words in a later free group is an equality in its free
monoid. Thus the word-length and cancellation arguments
there are valid here. Supports and the strong-singleton
tree-order argument are unchanged. Consequently either
both families are constant already, or all the same-family
and cross-family distinct-point values are pairwise
incomparable on both sides.

Assume the latter case. At one strong edge $a<b<c$, choose
a finite free chart witnessing (G.25). The words $u_a,v_c$
are not prefix-comparable. Remove their longest common
prefix $p$ in this chart, writing $u_a=pq$ and $v_c=ps$,
where $q,s$ have distinct first letters. In the reduced word
$u_bq^{-1}s v_a^{-1}$, all of $q^{-1}$ must cancel against
the end of $u_b$. Hence for a positive word $r$,
\[
 u_a=pq,\qquad u_b=rq,\qquad
 v_c=ps,\qquad v_a=rs.
 \tag{G.26}
\]
The factors belong to $S$ and have finite strong supports;
they need not be supported on the edge itself.

Consider the orbit $P=G\cdot p$ with prefix order. If $p$ is
$G$-fixed this is already a chain. Otherwise (Op) identifies
its stabilizer with $G_F$ for a nonempty finite strong $F$,
so its orbit is $\Omega_F$. Prefix order has chain principal
downsets by Levi's lemma in finite free charts. By
Theorem~\ref{G:alltreeorders}, its comparability components
are chains. Since $p$ prefixes $u_a$, every element of
$G_a\cdot p$ prefixes $u_a$; these elements are pairwise
comparable. Thus $G_a$ stabilizes the component of $p$.
Likewise $p$ prefixes $v_c$, so $G_c$ stabilizes that component.
The stabilizer-generation identity gives
$\langle G_a,G_c\rangle=G$, because $a\ne c$ are strong
singletons. Therefore $P$ has only one component and is a
global chain. Applying the same argument to $r$, using its
upper words $u_b$ and $v_a$, proves that $R=G\cdot r$ is
also a global chain.

There is a comparison between these two chains: choose
$h\in G$ with $h(b)=a$. Both $p$ and $hr$ prefix $u_a$.
We use a general tree-order observation. If two invariant
transitive orbits $P,R$ are chains and some $p\in P$ lies
below some $r\in R$, the set of $P$-elements lying below
some $R$-element is nonempty and invariant, hence equals $P$.
For arbitrary $p'\in P,r'\in R$, choose $r''\in R$ above
$p'$. Compare $r',r''$ in their chain. Either $p'\preceq r'$
by transitivity, or $p',r'$ both lie below $r''$ and are
comparable by the chain-downset property. Reversing the
roles handles the other initial orientation. Hence every
element of $P$ is comparable with every element of $R$.
In particular $p$ and $r$ are prefix-comparable.

If they are distinct, their prefix orientation is strict and
is preserved by every $g\in G$. Transport (G.26) to the six
strong edges of the eleven-vertex structure (G.10), choosing
an automorphism for each edge. All transported point words
are the prescribed $u_x,v_x$ by equivariance. Put these
finitely many factors, and their strict prefix quotients,
in one finite free chart. For $e=(a,b,c)$ its word lengths
satisfy
\[
 \tau_e=|u_b|-|u_a|=|r_e|-|p_e|
       =|v_a|-|v_c|.
\]
All six $\tau_e$ have one strict sign. The two exact cycle
identities (G.11) give
$\tau_1+\tau_2+\tau_4+\tau_5=0$, a contradiction.
Thus $p=r$. Equation (G.26) gives $u_a=u_b$ and $v_a=v_c$.
An equivariant strong-singleton family with equal values at
two distinct points is constant, by
$\langle G_a,G_b\rangle=G$ and the analogous identity for
$a,c$. This contradicts the nonconstant case and proves
the theorem.
\end{proof}

Theorem~\ref{G:noincidencelift} excludes the full DF incidence
group-relator route, rather than only a chosen word pattern.
It does not assert that every nonzero $T_1(M_{\ab})$ would
produce that relator. The arbitrary abelianization and
comparison questions remain separate gates.

\subsection{A finite obstruction to the natural-order word lift}

\begin{proposition}[Proved: the natural-order pullback bases do not lift to words]
\label{G:natural-order-fiber-obstruction}
For the independently ordered punctured $\mathbb F_2$-vector-space
category, let
\[
 N(B)=\mathbb Z^{\mathcal T(B)}
\]
be the natural-order function diagram from
Proposition~\ref{A:ordered-vector-integral-torsion}, with its
point-indicator bases. There is no free-monoid diagram with these
specified free bases whose abelianization maps are the pullback
maps of $N$. The obstruction already occurs on one finite
commuting diagram with objects of dimensions one, two, and three.
This assertion does not exclude a different choice of bases or a
stabilized positive realization of a different diagram.
\end{proposition}
\begin{proof}
Let $V$ be the punctured one-dimensional space, let $P$ be the
punctured two-dimensional space ordered $x_1<x_2<x_3$, and let
$Q$ be the punctured three-dimensional space with its binary labels
in the reference order $1<2<\cdots<7$. Define
\[
 a(P)=(1,2,3),\qquad b(P)=(1,4,5),
 \qquad v:V\to P,\quad v(V)=x_1.
\]
These are increasing linear embeddings and $av=bv$. The sets
of auxiliary natural orders have respective sizes $1,6,168$.
In the point-indicator bases the pullback of one order of $P$ is
the sum of all $28$ orders of $Q$ restricting to it.

Any word with this abelianization must list exactly those $28$
distinct generators, once each. Similarly the image of the sole
generator at $V$ in the free monoid at $P$ must list all six
generators once. Functoriality along $av=bv$ would therefore give
one word listing the $168$ generators at $Q$, in which every fiber
of either restriction $\mathcal T(Q)\to\mathcal T(P)$ is a
contiguous block. Thus the two six-part restriction partitions
would both be interval partitions of one linear order.

Two interval partitions with six nonempty blocks each have at
most $6+6-1=11$ nonempty pairwise intersections: their combined
ten interior cuts divide the line into at most eleven intervals.
Here all $36$ intersections are nonempty. In the six order types
\[
 012,\ 021,\ 102,\ 120,\ 201,\ 210
\]
on the locally indexed points of each plane, exact enumeration of
the $168$ ordered bases gives the joint fiber cardinalities
\[
 \begin{pmatrix}
 10&10&2&2&2&2\\
 10&10&2&2&2&2\\
 2&2&8&4&8&4\\
 2&2&4&8&4&8\\
 2&2&8&4&8&4\\
 2&2&4&8&4&8
 \end{pmatrix}.
\]
Every entry is positive and every row and column sums to $28$.
The attached exhaustive script constructs each natural order by
an ordered basis and verifies this matrix with integer arithmetic.
The interval-partition bound is contradicted, proving the claim.
\end{proof}

\subsection{A twisted lift beyond the excluded group relator}

\begin{example}[Proved: a twisted finite lift with both global cut chains]
\label{G:twisted-finite}
The positive pattern
\[
 u_a=pq,\qquad u_b=rq,\qquad v_c=sp,\qquad v_a=sr
 \tag{G.27}
\]
has the incidence relation
$[u_b]-[u_a]+[v_c]-[v_a]=0$ after abelianization,
but its four core words are freely independent in the free group
on $p,q,r,s$. There is a free-monoid diagram on the full
subcategory of $\Cdelta$ with objects $V,W,X_{11}$ from (G.10),
with ranks $2,6,1213$, in which every edge has form (G.27), all
transition homomorphisms are injective, the twenty-two point
words are pairwise bifix, and their only integral linear
relations are the six prescribed incidence relations.
Moreover, the six $q$-values form one strict suffix chain and
the six $s$-values one strict prefix chain, both ordered by the
signed lexicographic flag $(-a,-b)$ on an edge $(a,b,c)$.
The existence of an extension to all of $\Cdelta$ is not asserted.
\end{example}
\begin{proof}
The count relation is immediate from (G.27). Folding the four
positive loops $pq,rq,sp,sr$ yields three vertices and six
unoriented edges, hence a subgroup of rank four. More explicitly,
the vertices are $0,x,y$ and the positive edges are
\[
 0\xrightarrow{p,r}x,\quad x\xrightarrow q0,
 \quad0\xrightarrow s y,\quad y\xrightarrow{p,r}0.
\]
No further fold is possible. Choose the tree edges $0\xrightarrow p x$ and $0\xrightarrow s y$.
The remaining edges give the basis
$A=pq$, $B=rp^{-1}$, $C=sp$, $D=sr$; the four core words
are $A,BA,C,D$, which are again a free basis since
$B=(BA)A^{-1}$. The original DF
relator reduces to
$r p^{-1}s p r^{-1}s^{-1}$, which is nonempty.

We specify the larger finite example completely by labelled
word positions and identifications, without listing all 2898
positions separately. Write $U_i,V_i$ for initial formal words
with distinct position labels, of lengths
\[
\begin{split}
 (|U_i|)_{i=1}^{11}&=(124,126,133,129,130,132,139,134,131,128,134),\\
 (|V_i|)_{i=1}^{11}&=(134,140,131,133,129,135,127,129,132,133,135).
\end{split}
\]
Number all their positions consecutively, starting at $1$, in
the order $U_1,\ldots,U_{11},V_1,\ldots,V_{11}$. In the edge order
\[
 (1,2,9),\ (2,3,10),\ (3,4,11),\ (1,4,5),\ (5,6,7),\ (7,8,9),
\]
take the block lengths
\[
\begin{array}{c|rrrrrr}
 &0&1&2&3&4&5\\\hline
 |p|&18&58&84&37&100&128\\
 |q|&106&68&49&87&30&11\\
 |r|&20&65&80&42&102&123\\
 |s|&114&75&51&92&27&4
\end{array}
\]
For each edge, identify corresponding positions in the two
occurrences of each of $p,q,r,s$ in (G.27). In addition,
identify every $q$-block with the suffix of the $q$-block of
edge $0$ having the same length, and identify every $s$-block
with the prefix of the $s$-block of edge $0$ having the same
length. Take the equivalence relation generated by exactly
these identifications. Its $1213$ classes are the alphabet of
$M(X_{11})$; each $U_i,V_i$ becomes the word of its successive
classes. Put $M(V)=\langle u,v\rangle^*$ and
$M(W)=\langle p,q,r,s,d,e\rangle^*$, with point maps
\[
 (u_1,v_1)=(pq,sr),\quad (u_2,v_2)=(rq,e),
 \quad (u_3,v_3)=(d,sp).
\]
For an edge $(a,b,c)$, map $p,q,r,s$ to its indicated blocks,
$d$ to $U_c$, and $e$ to $V_b$. Map the point $i$ to
$(U_i,V_i)$. The defining identifications verify all composition
identities. These are all arrows and compositions in the stated
finite full subcategory.

Here are compact exact rank certificates. Give an alphabet
class its least position label. The point-count matrix has a
$16\times16$ minor of determinant $1$, with rows
\[
 1,2,3,4,5,6,7,9,10,11,12,13,14,15,17,19
\]
and columns
\[
 1,19,38,57,76,95,114,125,251,513,1048,1179,1307,1846,2108,2370.
\]
The six edge matrices have nonzero unit minors with respective
column sets
\[
\begin{array}{c|l|r}
0&1,19,125,251,384,1048&1\\
1&19,38,57,251,643,1179&1\\
2&38,57,76,643,1307,1846&1\\
3&1,38,95,125,643,1846&1\\
4&1,19,95,513,775,2108&-1\\
5&1,114,513,1048,1441,2370&-1
\end{array}
\]
where the last column is the determinant. Thus each edge
map has full abelian rank six; the defect theorem makes its
six images a code, so the monoid map is injective. The point
maps are injective as well. The relation matrix $D$, with
rows $u_b-u_a+v_c-v_a$, has determinant $-1$ on columns
$1,2,3,5,7,13$, so its image is primitive of rank six.
It lies in the kernel of the point-count map, whose rank is
sixteen. The primitive rank-six subgroup is therefore the
whole integral kernel.

The two exact verification scripts supplied with this example
reconstruct and check every word composition, both global chains,
pairwise bifixness, the class count, and all the displayed minors
with integer arithmetic; their saved data include every full word.
The chain order of the six edges is $5,4,2,1,3,0$, which is
exactly increasing $(-a,-b)$. This proves all the finite claims.
\end{proof}

For any global realization of (G.27) in a smooth locally free
$G$-monoid, Theorem~\ref{G:alltreeorders} and (St) force the
$q$-orbit to be one suffix chain: its $G_a$ and $G_b$ orbits
are suffixes of the fixed words $u_a,u_b$, so both stabilizers
preserve its comparability component, and together generate
$G$. The $s$-orbit is similarly one prefix chain, using
$v_a,v_c$. Example~\ref{G:twisted-finite} shows that imposing
these two necessary global-chain conditions does not by itself
exclude even the full $X_{11}$ configuration. The mixed roles
of $p,r$, prefix in a $u$-word and suffix in a $v$-word, prevent
us from applying Theorem~\ref{G:noincidencelift}. Compatibility
under all further strong embeddings remains the unproved step.

\subsection{The remaining universal question}

The gap theorem settles the first prescribed test, including the
full category's non-strong raw points.  Propositions
\ref{G:saturation} and \ref{G:divisibility} reduce a possible
counterexample to a finite-rank injective diagram whose smooth
locally free monoid admits a proper factor in the same $G$-orbit.
Proposition~\ref{G:idempotent} further forces its ordinary
monoid colimit to contain a nonidentity idempotent, with the
two absorption relations (G.7).
Proposition~\ref{G:boundedrank} excludes every uniform finite
rank bound, and forces infinite rational rank even after the
filtered colimit over inclusions.
The gap example has precisely such proper-factor phenomena, yet
passes; that necessary condition is not a sufficient obstruction.
Its ordinary colimit is the monoid in presentation (G.5), with
the nonidentity idempotent $e$.  It is its classifying space,
not that monoid, which is equivalent to the classifying space
of the free group in (G.3).
Theorem~\ref{G:embeddinggaps} also excludes every gap diagram
on a lexicographic embedding action of a fixed strong source,
even though its ordinary colimit can have arbitrarily long
finite chains of nonidentity idempotents.
For the general class $\mathcal F$, Proposition
\ref{G:generalcomparison} excludes rank-one separators over
all total indices, while Proposition~\ref{G:boundedcoefficients}
forces any abelianization obstruction on such an index to
escape every global finite rank bound. Proposition~\ref{G:boundedtarget}
also rules out absorption in the bounded-rank colimit.
Theorem~\ref{G:groupcriterion} replaces the remaining comparison
and homology questions by the intrinsic smooth locally free
$G$-monoid tests (G.19)--(G.20), and proves their independence
of the total Fra\"iss\'e presentation. It does not establish
that all such tests pass. Theorem~\ref{G:alltreeorders}
classifies every invariant tree order on a finite-support orbit.
Its application in Theorem~\ref{G:noincidencelift} excludes the
entire original DF incidence group-relator construction, including
all rectangle factorizations and extra positive generators. It is
not proved that an arbitrary nonzero derived abelianization class
would yield that particular relator. Example~\ref{G:twisted-finite}
keeps a different positive lift of the abelian incidence relation
alive on the complete finite $V,W,X_{11}$ test, including the two
necessary global cut chains. Its extension under all further strong
embeddings remains unproved.

At the present checkpoint we have not proved or refuted the
existence of $M:A\to\FreeMon$ with
$L_1\colim_A M_{\ab}\ne0$, nor the universal comparison assertion
for $A$, nor strictness of either inclusion
$\mathrm{exact}\subseteq\mathcal F\subseteq\mathrm{total}$.
The exact outstanding assertions and attempted constructions are
kept in \texttt{notes/G\_progress.md}; none of these three assertions
is asserted as a theorem here.

\paragraph{Literature checked.}
Fiedorowicz's opened Proposition 4.4 identifies asphericity of a
monoid classifying space with the vanishing of the positive Tor
groups of its group-completion module \cite{GFiedorowicz1984}.
This is a criterion for the target, not a proof of our comparison.
Gray--Steinberg, Remark 4.11, gives the explicit counterexample to
Fiedorowicz's flatness Lemma 5.2(a) \cite{GGraySteinberg2024}; we
have not used that lemma or its dependent amalgamation argument.
Holub's revised Lemma 3 proves freeness of the coincidence monoid
of nonerasing word morphisms \cite{GHolub2003}.  Its common-boundary
factorization checks a local overlap, but does not supply the
higher coherences needed for the universal assertion.
Searches for ``exact colimits free monoids'', ``totally aspherical
free monoids'', ``free monoids homotopy colimit'', and ``locally
free monoids homology'' did not locate the gap theorem or settle
the remaining assertion.  This is a search report and carries
no claim of priority.

\subsection{Positive abelian antichains and synchronized twisted cuts}

\begin{lemma}[Proved: all invariant orders on strong singletons]
\label{G:singletonpartial}
Every $G$-invariant partial order on the strong-singleton orbit
$\Omega$ is equality, the raw linear order, or its reverse.
No assumption on principal downsets is required.
\end{lemma}
\begin{proof}
Suppose $x\prec y$ with $x\ne y$, and choose a finite strong
$D$ containing them. First suppose $x<y$ in the raw order.
Form a path of four free copies $D_1,\ldots,D_4$, identifying
$y$ in $D_i$ with $x$ in $D_{i+1}$. Denote the five distinguished
vertices by $x_0,\ldots,x_4$. Each gluing is over a strong
singleton. The raw orders extend with
$x_0<x_1<\cdots<x_4$, and invariance and transitivity give
$x_0\prec x_4$.

We verify that the endpoint pair is strong. For a subset $Y$
containing both endpoints, if all three internal gluing vertices
belong to $Y$, each $Y\cap D_i$ contains two distinct strong
singletons, so $\delta(Y\cap D_i)\ge3$. Consequently
$\delta(Y)\ge4\cdot3-3\cdot2=6>4$.
If an internal gluing vertex is missing, cut the path at every
such vertex. The two endpoint components each have predimension
at least $2$, with equality only when that component is its bare
endpoint: the first nontrivial extension of a strong singleton
has strictly positive increment, and subsequent free pieces have
nonnegative increment. Every nonempty middle component has
positive predimension. Thus every strict enlargement of the
endpoint pair has predimension greater than $4$.
The pair is therefore a strong copy of the ordered two-point
structure with no edge. Homogeneity now gives $a\prec b$ for
every free strong pair $a<b$.

For arbitrary strong singleton points $a<b$, choose a finite
strong support containing them and freely adjoin an isolated
point $z$ between them. Both $\{a,z\}$ and $\{z,b\}$ are strong
free pairs, so $a\prec z\prec b$. This proves the raw linear
order. The case $y<x$ gives its reverse by the same construction.
If there was no distinct comparable pair, the order was equality.
\end{proof}

\begin{proposition}[Proved: abelian antichains synchronize the twisted cuts]
\label{G:twisted-synchronization}
Let $S$ be a smooth locally free $G$-monoid, and let
$x\mapsto u_x$ and $x\mapsto v_x$ be equivariant maps from
$\Omega$ to $S$. Suppose that in $S_{\ab}$
\[
 [u_b]-[u_a]+[v_c]-[v_a]=0
\]
for every strong ordered edge $(a,b,c)$.
Then either both abelianized point families are constant, or
both are injective and are antichains for the pointed cone
$P=\operatorname{im}(S\to S_{\ab})$.

Assume the nonconstant case. If equivariant edge factors
$p_e,q_e,r_e,s_e\in S$, for $e=(a,b,c)$, satisfy
\[
 u_a=p_eq_e,\qquad u_b=r_eq_e,\qquad
 v_c=s_ep_e,\qquad v_a=s_er_e,
\]
then each of $q_e,s_e$ has full strong support $e$.
Their respective suffix and prefix orders are global linear
orders on the edge orbit, and they have exactly the same
signed two-coordinate lexicographic flag. Thus only the
$24$ identical-flag regimes remain possible, among the $576$
pairs of flags.
\end{proposition}
\begin{proof}
The cone $P$ is closed under addition, is pointed, and the
abelianization of a nonidentity element of $S$ is nonzero.
Indeed, every finite assertion is witnessed in a finite free
chart, and passage to a larger chart sends each nonempty word
to a nonempty word. A nonzero vector with nonnegative coordinates
cannot become zero, or become the negative of another such
vector, at a later chart.

For either point family, equality at two different points makes
its common value fixed by both singleton stabilizers. By (St)
these generate $G$, so the whole family is constant. If one
abelianized family is constant, the edge identity gives equality
of the other family at two distinct points, so both are constant.
Otherwise both maps to $S_{\ab}$ are injective. Pulling back the
order defined by $P$ gives invariant partial orders on $\Omega$.
By Lemma~\ref{G:singletonpartial}, a single strict comparison
would make the corresponding family strictly increasing or
strictly decreasing in raw order. On the six-edge object
(G.10), write $\tau_e=[u_b]-[u_a]=[v_a]-[v_c]$.
The integral identity (G.11) says
$\tau_1+\tau_2+\tau_4+\tau_5=0$.
Strict monotonicity of either family puts all four nonzero
summands in $P$, or all in $-P$, which contradicts pointedness.
Both families are therefore antichains.

The suffixes of any one word in $S$ are linearly ordered,
as are its prefixes; this follows in a common free chart.
The $G_a$ and $G_b$ orbits of $q_e$ lie among the suffixes
of $u_a$ and $u_b$, respectively. Each stabilizer preserves
the comparability component of $q_e$, and (St) makes that
component the whole $G$-orbit. Theorem~\ref{G:alltreeorders}
therefore makes this orbit one suffix chain. Likewise the
$G_a$ and $G_c$ orbits of $s_e$ lie among prefixes of
$v_a,v_c$, and its whole orbit is one prefix chain.
Set $Q_e=[q_e]$, $S_e=[s_e]$ and
\[
 H_e=S_e-Q_e=[v_c]-[u_a]=[v_a]-[u_b].
\]
If $e,f$ share their $a$ coordinate and their two chain orders
have opposite strict directions, or one is a tie and the other
is strict, then $H_e-H_f$ or its negative is a nonzero element
of $P$. It equals $[v_c]-[v_{c'}]$, contradicting the antichain
property. The same argument for a shared $c$ uses
$-[u_a]+[u_{a'}]$, and for a shared $b$ uses
$[v_a]-[v_{a'}]$. Thus the two chain comparisons, including
ties, agree on every fiber of each of the three coordinate maps.

By (Op), a stabilizer of an edge factor is the pointwise
stabilizer of a strong subset of its edge. The proper strong
subsets of a three-point edge are the empty set and its three
singletons. Suppose $q_e$ has such a proper support. Choose a
singleton coordinate $x$ whose stabilizer fixes $q_e$.
The fiber comparison just proved makes $s_e$ fixed by $G_x$
as well. If $x=a$, cancellation in $v_a=s_er_e$ makes $r_e$
fixed by $G_a$, and $u_b=r_eq_e$ is then fixed by both
$G_a,G_b$. If $x=b$, cancellation in $u_b=r_eq_e$ makes
$r_e$ fixed by $G_b$, and $v_a=s_er_e$ is fixed by both
$G_a,G_b$. If $x=c$, cancellation in $v_c=s_ep_e$ makes
$p_e$ fixed by $G_c$, and $u_a=p_eq_e$ is fixed by both
$G_a,G_c$. In every case (St) makes a point word globally
fixed, contradicting the nonconstant abelianized case.
Thus $q_e$ has full edge support. The same argument begins
with a proper support of $s_e$, since the fiber comparisons
include ties, so $s_e$ has full edge support too.

Theorem~\ref{G:alltreeorders} now identifies each chain with a
signed lexicographic order on two distinct coordinates of the
edge: the closure of one coordinate is its singleton and the
closure of any two coordinates is the whole edge. These flags
are determined by their restrictions to the three
single-coordinate fibers. Indeed, if their first pivots differ,
fix the third coordinate and place two free edge copies over
that singleton so that the other two coordinate comparisons
make the proposed first-pivot decisions disagree. If their
first pivots agree but their signs differ, fix any other
coordinate and vary the first. Finally, if the first pivot and
sign agree, fix that coordinate; arbitrary sign patterns on
the two remaining coordinates distinguish different second
pivots or different second signs. Each such placement exists
by ordered free amalgamation over the fixed singleton, with
small independent coordinate intervals. The fiber agreements
force equality of both flags, as claimed.
\end{proof}

This excludes different global cut flags for the twisted
incidence construction, without excluding its remaining
identical-flag regimes. Proposition~\ref{G:twisted-finite}
provides an exact finite model in one such regime. No global
extension, nonzero derived class, or unrestricted comparison
conclusion is asserted here.

\subsection{Canonical ordered cuts and their stabilizers}

\begin{proposition}[Proved: bounded cut stabilizers are locally cyclic]
\label{G:cutstabilizers}
Let $S$ be any locally free monoid, of arbitrary cardinality, and
put $K=S^{\gp}$. The map $S\to K$ is injective, $K$ is locally
free as a group, and
\[
 x\le y\quad\Longleftrightarrow\quad x^{-1}y\in S
 \qquad(x,y\in K)
\]
is a left-invariant partial order with linearly ordered intervals
and the interpolation property
\[
 x_1,x_2\le y_1,y_2
 \quad\Longrightarrow\quad
 \exists k\in K\quad x_1,x_2\le k\le y_1,y_2.
\]
For $E\subseteq K$, write $E^u$ and $E^l$ for its common upper
and lower bounds. A bounded Dedekind--MacNeille cut is a pair
$c=(D,U)$ satisfying
\[
 D=U^l\ne\varnothing,\qquad U=D^u\ne\varnothing.
\]
Its stabilizer $K_c=\{g\in K:gD=D\}$ under left translation
maps injectively to $K_{\ab}=S_{\ab}$ and is locally cyclic.

If a topological group $G$ acts smoothly on $S$ by monoid
automorphisms, its induced action on the cuts obtained from
principal cuts by finitely many bounded meets and joins is
smooth. For $G=\Aut(M_0)$ of $\Cdelta$, let $\mathscr D$ be
any invariant smooth set of bounded cuts. Each cut stabilizer
$G_c$ fixes $K_c$ pointwise. Consequently, for every unital ring
$R$, including the zero ring, the coefficient diagram
\[
 N_R(B)=\bigoplus_{c\in\mathscr D^{G_B}}
            R\otimes_{\Z}K_c
\]
has $L_i\colim_{\Cdelta}N_R=0$ for every $i>0$.
The transition along an embedding extended by $g\in G$ sends
the summand at $c$ to the summand at $gc$ by $1\otimes g$.
These assertions do not assert that this cut completion is a
tree or supplies an acyclic interval resolution.
\end{proposition}
\begin{proof}
Finite free submonoids containing prescribed finite subsets
form a directed system with union $S$. Group completion
preserves its colimit. A finite positive equality in $K$ is
therefore witnessed in the group completion of a later free
chart, where positive words embed; hence $S\to K$ is injective.
The same finite-chart argument proves cancellation and conicality
in $S$, so the displayed relation is a partial order.
After translating an interval $[a,b]$ by $a^{-1}$, its points
are prefixes in $S$ of $a^{-1}b$. Levi's lemma in a common
free chart makes these prefixes linearly ordered.

We give the group-freeness check without assuming that chart
inclusions induce injective homomorphisms of free groups.
Fix finitely many positive elements of $S$. In each later free
chart, the subgroup they generate is finitely generated free,
by the covering-graph proof of the subgroup theorem. Its rank
is bounded by the size of the tuple. Maps between these image
subgroups are surjective, so after their rank attains its minimum
they are surjections between free groups of equal finite rank.
Such maps are isomorphisms. For completeness, a finitely
generated residually finite group is Hopfian: if a surjection
$\phi$ killed $g\ne1$, choose a finite-index normal subgroup
$N$ avoiding $g$. There are finitely many subgroups of that
index, so two subgroups $\phi^{-i}N$ coincide; surjectivity then
gives $N=\phi^{-d}N$ for some $d>0$, contrary to
$g\in\ker\phi\subseteq\phi^{-d}N$. A free group is residually
finite directly: realize a nontrivial reduced word as a path
on finitely many distinct vertices, extend the partial bijections
specified by each generator along that path to permutations,
and the resulting finite permutation image does not kill the
word. Thus the image subgroup of our tuple in $K$ is free.
Every finite tuple of elements of $K$ lies in a subgroup generated
by finitely many positive elements, proving local freeness.

For interpolation, represent the four vertices and all four
positive increments $x_i^{-1}y_j$ in a free chart and enlarge
it until their finitely many group equalities hold. In the
oriented Cayley tree of that free group, the four directed paths
from the two lower vertices to the two upper vertices have a
common intermediate vertex: their incoming portions merge
before their outgoing portions split. Its image in $K$ is the
required $k$.

We will also use positive nonvanishing:
\[
 s\in S,\quad [s]=0\text{ in }S_{\ab}
 \quad\Longrightarrow\quad s=1.
\]
Indeed $S_{\ab}$ is the filtered colimit of the free abelian
chart groups. A nonempty positive word has a nonzero vector
with nonnegative coordinates, and every subsequent chart
inclusion sends it to a nonempty word. Its class cannot vanish.

Fix a bounded cut $(D,U)$. We first prove a cut dichotomy.
If $a\in D$, $b\in U$ and $a\le k\le b$, then
$k\in D\cup U$. If $k\notin D$, choose $u'\in U$ such that
$k\not\le u'$. For any $z\in D$, interpolation supplies
$h$ with $a,z\le h\le b,u'$. The elements $h,k$ lie in the
linearly ordered interval $[a,b]$. The inequality $k\le h$
would imply $k\le u'$, so $h<k$ and hence $z\le k$.
Thus $k$ is above all of $D$, which means $k\in U$.

It follows that either $D$ is upward directed or $U$ is
downward directed. If $D$ is not directed, choose $x,y\in D$
with no common upper bound in $D$. For any $u,v\in U$,
interpolation gives $x,y\le k\le u,v$. The cut dichotomy
places $k$ in $D\cup U$, and the choice of $x,y$ excludes
$D$. Thus $k\in U$ is a lower bound of $u,v$ within $U$.

Let $g\in K_c$, which also preserves $U$. If $D$ is directed,
fix $a\in D,b\in U$ and choose $x\in D$ with
$x\ge a,g^{-1}a$. Then $x,gx\in[a,b]$, so they are comparable.
If instead $U$ is downward directed, choose $x\in U$ with
$x\le b,g^{-1}b$, and the same conclusion holds.
Consequently either $x^{-1}gx\in S$ or its inverse belongs
to $S$. If $g$ maps to zero in $K_{\ab}=S_{\ab}$, positive
nonvanishing forces $g=1$. This proves the asserted injection.
It makes $K_c$ abelian; every finitely generated subgroup is
also free, by local freeness of $K$, and therefore cyclic.

Finally, every group element of $K$ is represented by finitely
many positive elements, so the given action of $G$ on $K$ is
smooth. A principal cut is determined by one such element.
Meets of closed lower sets are intersections, and joins are
the common-lower-bound sets of their common upper bounds;
these operations commute with order automorphisms. A finite
term in bounded meets and joins is therefore fixed by the
intersection of the open stabilizers of its finitely many
parameters. This proves smoothness.

For the coefficient assertion, (Op) gives $G_c=G_F$ for a
finite strong $F$, possibly empty. If $K_c\ne1$, choose
$h\in K_c\setminus\{1\}$. An automorphism of the torsion-free
rank-one group $K_c$ that fixes $h$ fixes every element:
for each element some nonzero integral multiple is an integral
multiple of $h$. Hence the kernel of the $G_F$-action on $K_c$
is the open stabilizer of $h$, and it is normal in $G_F$.
There is no proper open normal subgroup of $G_F$. Indeed,
(Op) writes one as $G_D\subseteq G_F$, so (Fix) gives
$F\subseteq D$. Normality makes the finite set $D$ invariant
under $G_F$; the raw linear order makes this action pointwise
trivial, and (Fix) then gives $D\subseteq F$.
Thus $G_F$ fixes $K_c$ pointwise. The case $K_c=1$ is immediate.

It follows that the bundle of groups $K_c$ over each orbit of
cuts is the constant group $K_c$ transported over $G/G_F$.
The fixed-cut coefficient diagram for this orbit is
$(R\otimes_{\Z}K_c)\otimes_{\Z}\Z[\operatorname{Emb}(F,-)]$.
For $F\ne\varnothing$ its bar complex is the chain complex
of $F/\Cdelta$, with constant coefficient $R\otimes K_c$;
that category has an initial object. For $F=\varnothing$ it
is the constant-coefficient complex of $B\Cdelta$, which is
contractible by totality. Both have zero positive homology,
without any flatness assumption on $R$. Direct sums over cut
orbits prove the assertion. The pointwise stabilizer statement
also makes the displayed transition maps independent of the
chosen extension of the embedding.
\end{proof}

The locally cyclic conclusion cannot in general be strengthened
to trivial stabilizers, even for a cut defined by two elements.
Let $S$ be the sequential union of the free monoids
$\operatorname{Free}(a,b_n,c_n)$ under
\[
 a\longmapsto a,\qquad
 b_n\longmapsto a b_{n+1},\qquad
 c_n\longmapsto a c_{n+1}.
\]
The three images form a code, and the induced group maps
are isomorphisms; hence $K$ is the free group on $a,b_0,c_0$.
In the $n$th chart, $b_0=a^n b_n$ and $c_0=a^n c_n$.
Their common lower set in $K$ is
$D=\bigcup_{n\ge0}\mathord\downarrow a^n$: any proposed common
lower vertex and its two positive increments can be put in a
later free chart, whose greatest common lower vertex for
$b_0,c_0$ is $a^n$.
This lower set is the nonprincipal meet cut of $b_0,c_0$,
and $aD=D$. Its stabilizer is abelian by the proposition,
so it centralizes $a$. Reduced words in the free group show
that the centralizer of a basis letter is its cyclic subgroup:
a word containing a different letter cannot commute with $a$
after removing its initial and terminal powers of $a$.
The cut stabilizer is therefore exactly $\langle a\rangle\cong\Z$.

\subsection{A canonical tree completion of the prefix order}

\begin{proposition}[Proved: finite-parameter median tree completion]
\label{G:prefixtree}
Let $S$ be a locally free monoid, $K=S^{\gp}$, and $P=(K,\le)$
the prefix order of Proposition~\ref{G:cutstabilizers}. Inside its
bounded Dedekind--MacNeille completion, let $T$ be the closure of
the principal cuts under finitely many bounded binary meets and
joins. There is a canonical median pretree structure on $T$ whose
segment between comparable points is their full order interval.
Every segment is a finite concatenation of order intervals, and
its turns are obtained by bounded binary meets and joins.
The order-preserving actions of $K$ and of every group of monoid
automorphisms of $S$ preserve this pretree structure. If the latter
action is smooth, its action on $T$ is smooth. Every $K$-point
stabilizer in $T$ is locally cyclic. For every subgroup $U$ of the
automorphism group, the fixed-point poset $T^U$ has weakly
contractible nerve; no openness assumption on $U$ is needed.

This supplies a tree of ordered intervals, not a simplicial tree
with a permutation module of edges. No equivariant homology
resolution or free-monoid comparison follows here by assertion.
\end{proposition}
\begin{proof}
We give the completion and path arguments explicitly. The general
connection between ramification-complete diamond-free orders and
simply connected nerves is also recorded in
\cite[Theorem~3.3.7]{GChen2014}; the finite path argument below
establishes what is needed without a countability restriction.

\emph{Intervals in the completion are chains.}
Suppose two incomparable bounded cuts $B,C$ lay between bounded
cuts $A,D$. Choose original vertices $a\in\operatorname{lower}(A)$
and $d\in\operatorname{upper}(D)$. Since the cuts $B,C$ are
incomparable, choose
\[
 b\in\operatorname{lower}(B),\quad
 u_C\in\operatorname{upper}(C),\quad b\not\le u_C,
 \qquad
 c\in\operatorname{lower}(C),\quad
 u_B\in\operatorname{upper}(B),\quad c\not\le u_B.
\]
Interpolation in $P$ gives
\[
 a,b\le h_b\le d,u_B,\qquad
 a,c\le h_c\le d,u_C.
\]
Both $h_b,h_c$ lie in the original chain $[a,d]$.
The relation $h_b\le h_c$ would give $b\le u_C$, and the reverse
relation would give $c\le u_B$. This contradiction proves the
claim, in particular for $T$. Any two points of $T$ with a common
upper bound in $T$ have their join in $T$, and the dual statement
holds for meets, by its defining closure.

\emph{The nerve of $T$ is simply connected.}
For a finite free chart $S_0\subseteq S$, its free group has the
oriented Cayley tree and the corresponding prefix poset $P_0$.
The nerve of $P_0$ is weakly contractible. Indeed, the finite
Cayley-tree hull of any finite collection of its vertices is a
finite oriented tree. Its order complex contracts by successively
removing a leaf: every vertex comparable with that leaf is
comparable with its unique neighbouring vertex, which cones its
link. Every finite simplicial sphere therefore contracts after
passing to such a hull.

The posets $P_0$ need not embed in one another. Nevertheless their
nerves have filtered colimit $NP$: every finite chain in $P$, its
positive increments, and every equality between two such finite
records are witnessed in a later free chart. Consequently $NP$
is weakly contractible, by the finite simplicial description of
homotopy groups. In particular it is connected and simply
connected.

Adjoin the nonprincipal vertices of $T$ to $P$ in any well-order,
with the induced partial order at every stage. Each new cut has
an original lower and an original upper vertex. Its link in the
preceding order complex is
\[
 N(\{x:x<c\})*N(\{x:c<x\}),
\]
whose two factors are nonempty, so the link is connected.
Adjoining its cone cannot create a fundamental-group generator:
van Kampen makes the map on fundamental groups surjective.
The same remains true at limit stages, since an edge loop and
an edge-path homotopy use finitely many vertices. Thus $NT$ is
connected and simply connected.

\emph{Reduction of interval paths.}
An interval path is a finite list of nondegenerate order
intervals, successively traversed in either direction, with
matching endpoints. Consecutive traversals in the same direction
merge. At a peak $a\le b\ge c$, the two intervals overlap in
$[a\vee c,b]$; cancel that retraced portion and replace the turn
by $a\vee c$. At a trough use $b$'s replacement $a\wedge c$.
These are actual points of $T$, and the overlap is closed.

Define reduction by a stack. Append the next interval to an
already reduced stack and compare it with the terminal interval.
If the common retraced part consumes the old interval, pop it
and continue. If it consumes the new interval, trim and stop.
If it consumes neither, trim both at the overlap endpoint and
stop. Delete zero-length pieces and merge equal-direction runs.
There are only finitely many pops. Shortening a terminal interval
cannot spoil the preceding reduced turn: for example, if its
old lower endpoint is $a=u\wedge b$ and $a\le e\le b$, then
$u\wedge e=a$; the dual identity handles a peak. Thus the output
has consecutive intervals intersecting only in their common
endpoint. Merging equal-direction pieces preserves the preceding
reduced turn as well: if $u\wedge b=a<b\le d$, then
$u\wedge d=a$. Indeed $u\wedge d$ and $b$ lie in the chain
$[a,d]$; the possibility $b\le u\wedge d$ would imply
$b\le u$, and otherwise $u\wedge d\le b$ forces
$u\wedge d\le u\wedge b=a$. The dual argument applies
at a peak.

Here is the local calculation ensuring that this operation is
independent of grouping. For three alternating intervals
$a\le b\ge c\le d$, set
\[
 e=a\vee c,\qquad f=b\wedge d.
\]
Both lie in the chain $[c,b]$. If $e\le f$, then
$e\wedge d=e$ and $a\vee f=f$, and either order of cancellation
gives the single interval from $a$ to $d$. If $f\le e$, then
$e\wedge d=f$ and $a\vee f=e$, and both orders give
$a\to e\to f\to d$, after deleting degenerate pieces.
The dual calculation covers the opposite orientation; a repeated
orientation is just interval concatenation.

For clarity, this calculation applies in a longer context as
follows. Compare the two groupings at the terminal pair of the
current stack. A partial cancellation leaves the preceding stack
unchanged, and the calculation gives the same new terminal
pieces. A whole-interval cancellation removes at least one stack
entry; apply the induction hypothesis to the shorter stack and
the remaining interval. Induction on the number of entries
removed, followed by induction on the number of appended
intervals, proves associativity of reduced concatenation.
It also proves that inserting a subdivision of an interval, or
inserting an interval followed by its reverse, leaves reduction
unchanged. These are exactly the triangle and backtrack moves
for edge paths in the order complex. Reduction is therefore
invariant under every finite edge-path homotopy in $NT$.

\emph{Unique segments and medians.}
Since $NT$ is simply connected, a closed interval path reduces
to the empty path. Two reduced paths with the same endpoints
must therefore agree, after monotone subdivisions: concatenate
one with the reverse of the other and cancel their common
terminal portions. Any remaining portion would be a nonempty
reduced closed path. A reduced path has no nonadjacent
self-intersection, for such an intersection would similarly
produce a nonempty reduced closed subpath. Its image, ordered
along the traversal, is the canonical segment between its
endpoints.

The intersection of two order intervals, when nonempty, is the
closed interval between the join of their lower endpoints and
the meet of their upper endpoints. Hence the intersection of
two finite interval paths is a finite union of closed intervals.
Unique paths imply that two segments from the same initial point
share an initial segment with a last point. For three points
$x,y,z$, this last point of $[x,y]\cap[x,z]$ lies on $[y,z]$;
otherwise their paths would contain a reduced cycle. It is the
unique median. These unique segments and their tripod
intersections give the pretree betweenness axioms directly.

Everything was defined by the partial order and bounded meets
and joins. The indicated actions therefore preserve it.
Smoothness and local cyclicity are exactly the final assertions
of Proposition~\ref{G:cutstabilizers}.

Finally, let $U$ be any subgroup of the given automorphism group.
It fixes the original vertex $1$. The unique path between two
$U$-fixed vertices is preserved in its traversal order. Its
finitely many order turns are therefore individually fixed.
For finitely many fixed vertices, take the union of their paths
from $1$, and retain the endpoints, all turns, and all branch
points. This gives finitely many vertices: adjoining one root
path meets the preceding union in an initial segment with a last
point, and adds only its finitely many turns and that possible
branch point. Every retained vertex is fixed, since branch points
are medians of fixed vertices. The induced finite poset is the
order of an oriented finite tree: its edges are the successive
retained vertices along these paths, and comparability is exactly
the existence of a directed path. Its nerve contracts by the leaf
removal argument above. Any finite simplicial sphere in $NT^U$
is contained in such a contractible finite subposet. Hence
$NT^U$ is weakly contractible.
\end{proof}


\subsection{Fixed monoids and fixed group completions differ}

\begin{example}[Proved: neither direction of the fixed-group identification]
\label{G:fixedgroupfailures}
For smooth locally free monoids $S$ of the group $G$ of $\Cdelta$,
the canonical homomorphism
\[
 (S^G)^{\gp}\longrightarrow (S^{\gp})^G
\]
need be neither injective nor surjective. The two explicit
examples below come from finite-rank free-monoid diagrams over
$\Cdelta$. Both diagrams satisfy $T_i(M_{\ab})=0$ for every
$i>0$, so these examples do not settle the remaining general
$T_1$ or comparison question.
\end{example}
\begin{proof}
Let $\Omega=\operatorname{Emb}(V,M)$, identified with the
strong-singleton vertices of the limit and ordered by the raw
linear order. This is an infinite transitive smooth $G$-set.
First form the poset $D$ with moving points $\Omega$, two fixed
incomparable left points $L_0,L_1$, and two fixed incomparable
right points $R_0,R_1$, with
$L_i<x<R_j$ for every $x\in\Omega$ and every $i,j$.
Let $S$ be the monoid generated by the nonidentity arrows of
this poset, with the relation that a pair of composable arrows
is their composite. We write products in path order.

Every element has a unique normal form consisting of a sequence
of nonidentity arrows in which adjacent arrows are not composable.
Indeed, composition strictly reduces the number of arrows;
disjoint reductions commute, and the sole overlapping ambiguity
is three composable arrows, resolved by associativity. Thus the
reduction is terminating and confluent; see also
\cite[Lemmas~3.3--3.4]{GWehrung2018} for the general category
normal form. A finite nonempty set
$I\subset\Omega$ of size $n$ gives the free submonoid on the
$n+3$ edges of the oriented tree with two incoming leaves,
the chain $I$, and two outgoing leaves. Every arrow involving
only the four fixed endpoints can be expressed using any point
of $I$. These free submonoids are directed and exhaust $S$;
inclusions subdivide edges and are injective by the normal form.
The action is smooth since each normal form uses finitely many
moving endpoints.

An element fixed by $G$ has every arrow in its normal form fixed.
The only fixed nonidentity arrows are
$w_{ij}=(L_i\to R_j)$. They cannot be composed with one another,
so $S^G$ is the free monoid on these four letters. At any
$x\in\Omega$, put $a_i=(L_i\to x)$ and $b_j=(x\to R_j)$.
Then $w_{ij}=a_i b_j$, whence in $S^{\gp}$
\[
 w_{00}w_{10}^{-1}w_{11}w_{01}^{-1}=1.
\]
The displayed word is reduced and nontrivial in the free group
$(S^G)^{\gp}$. This proves failure of injectivity.

For failure of surjectivity, omit the right endpoints and use
the analogous monoid $S'$ on $L_0,L_1$ and $\Omega$.
Its nonempty finite charts are free of rank $n+1$.
Every nonidentity arrow has a moving endpoint, so the same
normal-form argument gives $(S')^G=\{1\}$. Nevertheless
\[
 k=(L_0\to x)(L_1\to x)^{-1}
\]
is independent of $x$: if $x<y$, both arrows append the same
arrow $x\to y$, which cancels in the group. Hence $k$ is fixed
by $G$. It is nontrivial: the homomorphism $S'\to\Z$
sending an arrow starting at $L_0$ to $1$ and every other
nonidentity arrow to $0$ respects composition and sends $k$
to $1$ after group completion. Thus the fixed-group map is not
surjective.

We also check the asserted scope at the diagram level. For a
finite strong $B$, put $E(B)=\operatorname{Emb}(V,B)$, with its
induced order, and apply the two constructions with $E(B)$ in
place of $\Omega$. If $E(B)$ is empty, the two-sided monoid is
free on the four $w_{ij}$ and the one-sided monoid is trivial;
otherwise their ranks are $|E(B)|+3$ and $|E(B)|+1$.
A strong embedding acts on the moving endpoints and fixes the
boundary endpoints, so these are functors to free monoids.
Their inclusion colimits are $S$ and $S'$ respectively.

For the two-sided construction there is a natural surjection
\[
 M_{\ab}\longrightarrow\Z[E]\oplus\Z^3.
\]
When $E(B)\ne\varnothing$, its inverse sends the point basis
$x$ to $[L_0\to x]$, and the three constant basis vectors to
$[w_{00}],[w_{01}],[w_{10}]$. To verify these form a basis,
the internal edge $x\to y$ equals the difference of the two
point vectors; the terminal right edges are $[w_{0j}]$ minus
the last point vector; and the other initial left edge is
$[w_{10}]-[w_{00}]$ plus the first point vector.
At an empty stage the map sends
$[w_{11}]$ to $[w_{01}]+[w_{10}]-[w_{00}]$.
Its kernel is supported only on stages with $E(B)=\varnothing$.
In the one-sided construction there is instead a natural
injection
\[
 M'_{\ab}\longrightarrow\Z[E]\oplus\Z,
\]
which is an isomorphism at nonempty stages: the point vector
$x$ maps back to $[L_0\to x]$, and the constant vector to
$[L_1\to x]-[L_0\to x]$, independent of $x$.
Its cokernel is supported only on the empty stages.

Both exceptional diagrams are locally zero, since every stage
has a common successor with $V$. Proposition~\ref{G:moduleinjectivization}
kills all their derived colimits. The representable diagram
$\Z[E]$ and the constant diagrams have zero positive derived
colimits, the latter because $B\Cdelta$ is contractible.
The two exact sequences therefore give the claimed vanishing.
The same fixed-group phenomena already occur for
$\operatorname{Aut}(\Q,<)$ with the elementary ordered-set
charts; we have not used a claim that pullback from an exact
index preserves the comparison.
\end{proof}


\subsection{Splitting of discrete central stabilizer extensions}

\begin{proposition}[Proved: uniqueness and existence of a global central splitting]
\label{G:centralextensions}
In setting \textup{(S)}, suppose the strong-embedding category is
total. Let $G=\Aut(M)$, let $U\leq G$ be open, and let
\[
 1\longrightarrow L\longrightarrow E\stackrel{p}{\longrightarrow}U
 \longrightarrow1
\]
be an extension of topological groups with discrete abelian kernel,
where $p$ is a quotient map. Suppose the extension is central. Then it has a unique continuous
homomorphic section over $U$.

For the group of $\Cdelta$, in particular, a point $c$ of the
tree $T$ in Proposition~\ref{G:prefixtree} has a stabilizer in
$K\rtimes G$ fitting into
\[
 1\longrightarrow K_c\longrightarrow (K\rtimes G)_c
 \longrightarrow G_{Kc}\longrightarrow1,
\]
and this is canonically a direct product of $K_c$ and a uniquely
specified lift of $G_{Kc}$. Here $G_{Kc}$ denotes the setwise
stabilizer of the $K$-orbit of $c$, not the point stabilizer $G_c$.
Let $\Gamma=K_G$ be the group of coinvariants and let
$\rho:K\rtimes G\to\Gamma$ send $(k,g)$ to the class of $k$.
The uniquely lifted copy $\widehat{G_{Kc}}$ is killed by $\rho$.
Consequently
\[
 \rho((K\rtimes G)_c)=\rho(K_c),\qquad
 \ker(\rho)\cap(K\rtimes G)_c
 =\ker(\rho|K_c)\times\widehat{G_{Kc}}.
\]
In particular the image stabilizer in $\Gamma$ is locally cyclic.
This splitting does not assert that its lifted subgroup is
$K$-conjugate to the standard copy of $G_{Kc}$.
\end{proposition}
\begin{proof}
The existence of a continuous homomorphic local section is
automatic here. The quotient homomorphism $p$ is open. Since
its kernel is discrete, choose an open identity neighborhood
$O\subset E$ with $O^{-1}O\cap L=\{1\}$; then $p|O$ is a
homeomorphism onto the open set $pO$. Choose an open identity
neighborhood $W\subset O$ with $WW\subset O$. Since $U$ has
an open-subgroup basis, choose an open subgroup $V\subset pW$.
The inverse local section $s:V\to W$ is a homomorphism:
$s(x)s(y)$ and $s(xy)$ both lie in $O$ and have the same image
under $p$, which is injective there.

By Proposition~\ref{A:open-subgroups}, the open group $U$
has a total presentation in \textup{(S)}. Fix such a
presentation, and write $U_B$ for the pointwise stabilizer
of a finite strong set $B$ in its limit. By
Proposition~\ref{B:no-open-quotients}, no open subgroup of $U$
has a proper open normal subgroup. Thus every continuous
homomorphism from an open subgroup of $U$ to a discrete group
is trivial. Neither (Op) nor a free-amalgamation hypothesis
is being added to this assertion.

Let $V\leq U$ be an open subgroup on which a local section
has been obtained. Retain those finite strong sets $B$ for
which $U_B$ is contained in a $U$-conjugate of $V$.
These form a nonempty $U$-invariant upward-closed cofinal
collection. Its full embedding category $\mathcal A$ is a
cosieve of the total presentation, and is total: a set action
on this cosieve extends by the empty set outside it, while
common successors give connectedness. Its inclusion poset
$I$ is directed.

For each retained $B$ the extension splits continuously over
$U_B$. Such a section $s_B$ is unique. Indeed, the ratio of two
sections of a central extension is a continuous homomorphism
$U_B\to L$, whose open normal kernel must be all of $U_B$.
Consequently restrictions of these sections agree, and conjugation
by any element $e\in E$ carries $s_B$ to $s_{p(e)B}$.

Define a category $\widetilde{\mathcal A}$ with the same objects
as $\mathcal A$. A lifted arrow from $B$ to $C$ is a right coset
$e\,s_B(U_B)$ with $p(e)B\subseteq C$; its underlying arrow is
$p(e)|B$. Composition is multiplication of representatives.
It is well defined: if $p(e)B\subseteq C$, conjugation takes the
restriction of $s_C$ on $U_C\subseteq p(e)U_Bp(e)^{-1}$ into the
corresponding restriction of $s_B$, by uniqueness. Every arrow
of $\mathcal A$ therefore has an $L$-torsor of lifts, and the
central $L$-actions commute with composition.

Choose one lift of each arrow, choosing the identity lift for
identities. The defect of multiplicativity is a normalized
$L$-valued $2$-cocycle on the nerve of $\mathcal A$.
Since $B\mathcal A$ is weakly contractible, its cohomology in
degree two with these constant coefficients is zero. Subtract
a $1$-cochain whose coboundary is this defect. The corrected
lifts compose exactly. Hence
\[
 \widetilde{\mathcal A}\cong\mathcal A\times\mathbf B L,
 \qquad
 \pi_1(B\widetilde{\mathcal A})\cong L,
\]
and the latter isomorphism restricts to the identity on the
loops supplied by the kernel $L$.

We spell out how this categorical splitting gives a group
splitting. The transport category $I\rtimes E$ has the same
objects and has one arrow labelled $e$ from $B$ to $C$ whenever
$p(e)B\subseteq C$. Its nerve has fundamental group $E$:
localizing the inclusion arrows identifies all objects, because
$I$ is directed, and leaves precisely multiplication of the
labels in $E$. Equivalently its nerve is the homotopy quotient
of the contractible nerve $NI$ by the discrete group $E$.
The functor to $\widetilde{\mathcal A}$ is surjective on every
arrow set and imposes exactly the relations making each loop
in $s_B(U_B)$ an identity. Thus, if $N$ is the normal subgroup
of $E$ generated by all $s_B(U_B)$, localization gives
\[
 E/N\cong\pi_1(B\widetilde{\mathcal A})\cong L,
\]
with the original $L\subseteq E$ mapping identically onto this
quotient. In particular $N\cap L=1$ and $E=LN$, so
$p|N:N\to U$ is a group isomorphism.

It is a topological isomorphism. The image of a continuous local
section over a sufficiently small open subgroup is open in $E$:
discreteness of the kernel gives an identity neighborhood on
which $p$ is injective, and continuity of the section identifies
its restriction there with the inverse of this local quotient
homeomorphism. Therefore $N$ contains an open subgroup of $E$,
and is itself open. The quotient map $p$ is open, so its
restriction to $N$ is open as well. Its inverse is the required
continuous section. Uniqueness follows once more from the
absence of a nontrivial continuous homomorphism $U\to L$.

For the application, $G_{Kc}$ contains the open group $G_c$, hence
is open. Its full inverse image in $K\rtimes G$ preserving $c$
projects onto it, with kernel $K_c$. The subgroup
$\{1\}\times G_c$ gives a continuous local section.
By Proposition~\ref{G:cutstabilizers}, $K_c$ is locally cyclic
and $G_c$ fixes it pointwise. Conjugation on this abelian kernel
therefore descends to $G_{Kc}$, and its kernel contains $G_c$.
It is an open normal subgroup, so is all of $G_{Kc}$. The
extension is central, and the result just proved applies.
The projection is a quotient map because the displayed local
section has open image and its translates cover the stabilizer.
Finally, $\rho$ is a continuous homomorphism when $\Gamma$ is
discrete: its kernel contains the open subgroup $\{1\}\times G$.
Its composite with the unique lifted section over $G_{Kc}$ is
therefore a continuous homomorphism to a discrete group. Its
kernel is open and normal, so the composite is trivial. The two
displayed identities follow from the direct-product splitting.
An image of a locally cyclic group is locally cyclic; no
assertion of torsion-freeness of that image is made.
\end{proof}


\subsection{Affine fixed cuts do not imply conjugacy}

\begin{example}[Proved: a bounded affine fixed cut without a fixed group vertex]
\label{G:affinecutcounterexample}
There is a smooth locally free $G$-monoid $S$, for the group of
$\Cdelta$, and a continuous $1$-cocycle $\kappa:G\to S^{\gp}$
whose affine action on the tree of Proposition~\ref{G:prefixtree}
fixes a cut obtained by one bounded binary join, but fixes no
original group vertex. Thus the lifted subgroup in
Proposition~\ref{G:centralextensions} need not be conjugate by
$S^{\gp}$ to a subgroup of the standard copy of $G$.
\end{example}
\begin{proof}
Let $L$ and $\Omega$ be two tagged copies of the strong-singleton
$G$-orbit. Make $L$ an antichain, order $\Omega$ by the raw
linear order, and put every point of $L$ below every point of
$\Omega$. Denote this poset by $D$. Let $S$ be its interval
monoid, with the reduced-arrow normal forms used in
Example~\ref{G:fixedgroupfailures}. A chart on finite subsets
$L_0\subset L$ and $I\subset\Omega$, with $I$ nonempty, is the
free monoid on the oriented tree with one edge from each point
of $L_0$ to $\min I$ and the successive edges of $I$.
These charts inject and exhaust $S$, so it is locally free;
finite endpoints give smoothness.

Fix $o\in\Omega$. For $d\in D$, let $j(d)\in K=S^{\gp}$ be
the signed interval path from $o$ to $d$. Such a path is unique
in every finite chart containing $o,d$, and the definitions
agree under enlargement. In particular $j(o)=1$ and
$j(d)^{-1}j(e)=[d,e]$ whenever $d<e$ in $D$.
The group maps between charts containing $o$ are injective:
in vertex coordinates they are the inclusions of free groups
on the letters $j(d)$ for $d\ne o$. Thus $j$ embeds $D$ into
the prefix poset of $K$. This is an order embedding: in a
common oriented star-and-chain chart, the path between two
vertices is positive exactly for a comparable pair in $D$;
enlarging the chart does not change that criterion.

The continuous cocycle
\[
 \kappa(g)=j(go),\qquad
 \kappa(gh)=\kappa(g)\,g\kappa(h)
\]
gives the affine action $g*x=\kappa(g)g(x)$. It sends
$j(d)$ to $j(gd)$ and preserves the prefix order. Its
continuity follows from the open stabilizer of $o$.
The boundary homomorphism identifies, equivariantly,
\[
 K_{\ab}\cong\ker\bigl(\Z[D]\stackrel{\aug}{\longrightarrow}\Z\bigr),
 \qquad [d,e]\longmapsto e-d.
\]
Indeed this is the vertex basis calculation in the finite tree
charts, or follows from connectedness and vanishing first
homology of the directed poset $D$.
If $x\in K$ were fixed by the affine action, its abelian vector
$v$ would satisfy $v=go-o+gv$. Thus $v+o$ would be a
$G$-invariant finitely supported vector in $\Z[D]$ of augmentation
one. Both $G$-orbits in $D$ are infinite, so every finitely
supported invariant vector is zero. This is impossible.

It remains to verify the bounded cut in the full prefix order,
not merely in $D$. Choose distinct $l_1,l_2\in L$. Their
common-upper-bound set in $K$ is exactly
\[
 U=\bigcup_{y\in\Omega}\mathord\uparrow j(y).
\]
One inclusion follows from $j(l_i)\le j(y)$. For the other,
put any proposed common upper bound, its two positive increments,
and $o,l_1,l_2$ in one finite free chart, enlarging until the
finite equalities hold. In its oriented Cayley tree the join
of $j(l_1),j(l_2)$ is $j(\min I)$, where $I$ is that chart's
finite nonempty set of $\Omega$-vertices. Hence the proposed
upper bound lies above some $j(y)$, as required.
Consequently their bounded join cut $c$ has upper set $U$.
It is independent of the chosen distinct pair in $L$, is
bounded below by $j(l_1)$ and above by every $j(y)$, and is
fixed by the affine $G$-action. It belongs to the finite-parameter
tree because it is a single bounded binary join of original
vertices. It is nonprincipal, since a principal fixed cut would
be a fixed original vertex, already excluded.

The subgroup $\{(\kappa(g),g):g\in G\}$ of $K\rtimes G$
therefore stabilizes $c$. If it were conjugate into the
standard copy of $G$, its action on the $K$-torsor would fix
a vertex. It does not. The uniqueness assertion of
Proposition~\ref{G:centralextensions} identifies this subgroup
with the splitting over $G_{Kc}=G$.
\end{proof}


\subsection{Smooth quotients and the tree double complex}

\begin{lemma}[Proved: acyclic smooth spaces over a total group]
\label{G:smoothquotient}
Suppose a topological group $U$ has a total presentation in setting
\textup{(S)}. Let $X$ be a nonempty simplicial $U$-set whose action
in every degree is smooth. If $X$ is connected and integrally
acyclic, then $X/U$ is connected and integrally acyclic. If $X$ is
weakly contractible, then $X/U$ is weakly contractible.
Here the quotient is the degreewise simplicial quotient followed
by realization; no separation assumption on a topological orbit
space is being substituted for this construction.
\end{lemma}
\begin{proof}
Use the total strong presentation $\mathcal A$ and write $U_B$ for
its finite strong pointwise stabilizers. The simplicial sets
$X^{U_B}$ form a covariant diagram on $\mathcal A$: an embedding
acts by any automorphism extending it. Its action on a fixed
simplex is independent of that extension. The filtered colimit
over inclusions is $X$, since every finite collection of simplices
has a finite strong support.

Take the first-quadrant double complex given by the category bar
resolution of the unnormalized chain diagrams $C_q(X^{U_B};\Z)$.
Each chain diagram is the free abelian group on a set action.
Totality therefore makes its positive derived colimits zero:
they are the positive homology groups of the corresponding
categories of elements, whose components are contractible.
Its ordinary colimit is $\Z[X_q/U]$. One spectral sequence thus
identifies the homology of the double complex with $H_*(X/U;\Z)$.

For the other spectral sequence, every positive homology diagram
$B\mapsto H_q(X^{U_B};\Z)$ is locally zero. Indeed, a finite cycle
becomes a boundary in $X$, and a finite filling has a finite strong
support containing its original support. In degree zero, the
kernel and cokernel of the augmentation to the constant diagram
$\Z$ are locally zero, by the same finite-path argument and
nonemptiness. Proposition~\ref{G:moduleinjectivization} kills all
derived colimits of these locally zero diagrams. The other
spectral sequence consequently has just $\Z$ in total degree
zero, since the nerve of the total category is contractible.
This proves the homology assertion.

For completeness, if $X$ is simply connected, an elementary
edge-path argument gives a surjection
\[
 U\longrightarrow\pi_1(|X/U|).
\]
Choose a vertex $x$ and, for $u\in U$, an edge path from $x$ to
$ux$. Its projection is a loop. Simple connectivity makes its
class independent of the chosen path and makes this assignment
a homomorphism (up to the harmless opposite convention).
Every quotient edge loop can be lifted edge by edge, translating
successive lifts to match their endpoints; its final endpoint is
$ux$, proving surjectivity. An element fixing a vertex $y$ lies
in the kernel: use a path from $x$ to $y$ followed by the reverse
of its translate. The kernel therefore contains an open vertex
stabilizer. It is an open normal subgroup of $U$, and must be
all of $U$ by Proposition~\ref{B:no-open-quotients}. Thus the
quotient is simply connected. The homology assertion and the
Hurewicz theorem now make it weakly contractible.
\end{proof}

\begin{proposition}[Proved: a coherent tree model with locally cyclic fibres]
\label{G:treedoublemodel}
Let $G$ be the group of $\Cdelta$, let $S$ be a smooth locally
free $G$-monoid, and put $K=S^{\gp}$ and $H=K\rtimes G$.
Let $P$ be the prefix poset of $K$ and $T$ its tree completion.
There is a natural weak equivalence
\[
 \bigl|(NP\times NT)/H\bigr|\ \longrightarrow\ |(NS)/G|.
 \tag{TD.1}
\]
For a simplex $\sigma$ of $NT$, write $H_\sigma$ for its
pointwise stabilizer and $A_\sigma=K\cap H_\sigma$. Then
$A_\sigma$ is torsion-free locally cyclic and
\[
 |NP/H_\sigma|\simeq K(A_\sigma,1).
 \tag{TD.2}
\]
These fibres are the rows of one explicit bisimplicial set,
so all their incidence maps are retained. In particular there
is a homology spectral sequence converging to
$H_*((NS)/G;\Z)$, with
\[
 E^1_{n,q}=\bigoplus_{[\sigma]\in (NT_n)/H}
                H_q(A_\sigma;\Z),\qquad
 E^1_{n,q}=0\quad(q\geq2).
 \tag{TD.3}
\]
The $q=0$ row is the simplicial chain complex of $NT/H$; the
$q=1$ row has the indicated stabilizer groups and the face maps
induced by inclusions and conjugation. Formula (TD.3) alone does
not assert that either horizontal row is acyclic, or that the
remaining differential between them vanishes.
\end{proposition}
\begin{proof}
Use the bisimplicial set
\[
 Z_{m,n}=(NP_m\times NT_n)/H.
\]
Its diagonal is the quotient in (TD.1), and diagonal and iterated
realization have the same weak homotopy type.
For a fixed $P$-simplex $\tau$, its stabilizer in $H$ is conjugate
to an open subgroup of the standard $G$. Indeed, translate its
first vertex to $1$; fixing the remaining finitely many group
vertices is precisely fixing finitely many smooth group elements.
The stabilizer is open in $G$ and has a total presentation by
Proposition~\ref{A:open-subgroups}. Its action on $NT$ is smooth,
and $NT$ is weakly contractible by
Proposition~\ref{G:prefixtree}. Lemma~\ref{G:smoothquotient}
therefore makes $NT/H_\tau$ weakly contractible.
For every fixed $m$ the projection of the vertical realization
of $Z_{m,*}$ to $NP_m/H$ is consequently a weak equivalence.
Realizing in $m$ proves (TD.1), because
\[
 NP/H=(NP/K)/G=(NS)/G:
\]
a prefix chain modulo left translation records exactly its
successive positive increments in $S$.

Now fix a $T$-simplex $\sigma$. Its stabilizer projects onto
an open subgroup $U_\sigma\leq G$, since it contains the
standard open pointwise stabilizer of its finitely many smooth
vertices. Its kernel $A_\sigma$ is a subgroup of any one of
their locally cyclic $K$-stabilizers. The standard subgroup
fixing all vertices of $\sigma$ acts trivially on this kernel,
by Proposition~\ref{G:cutstabilizers}. Thus the conjugation
action of $U_\sigma$ on $A_\sigma$ has an open normal kernel
and is trivial. Proposition~\ref{G:centralextensions} gives
the unique continuous splitting
\[
 H_\sigma=A_\sigma\times\widehat U_\sigma.
 \tag{TD.4}
\]
The action of $\widehat U_\sigma$ on $NP$ is smooth. Since
$NP$ is weakly contractible, Lemma~\ref{G:smoothquotient}
makes $NP/\widehat U_\sigma$ weakly contractible.

The remaining $A_\sigma$-action on this quotient is free.
To see this, a stabilizer in $H_\sigma$ of an original vertex
$k\in K$ has trivial intersection with $A_\sigma$, because
$K$ acts freely on its original vertices. Its projection to
$U_\sigma$ is an open subgroup $V$, and it gives a continuous
section over $V$ of the central extension (TD.4). Uniqueness
in Proposition~\ref{G:centralextensions}, applied to this open
total group $V$, makes that section precisely
$\widehat U_\sigma|V$. Thus no nonidentity $A_\sigma$-element
can carry an original vertex into its $\widehat U_\sigma$-orbit.
The same follows for every simplex by considering a vertex.
The quotient of the weakly contractible simplicial set
$NP/\widehat U_\sigma$ by this free simplicial action is a
classifying space for $A_\sigma$, proving (TD.2).

Realize $Z_{*,n}$ first and take its homology. It is the disjoint
union of the spaces (TD.2), one for each orbit of $T$-simplices.
The standard first-quadrant spectral sequence of this bisimplicial
chain complex gives (TD.3). A torsion-free locally cyclic group
is the filtered union of its infinite cyclic subgroups, so its
integral homology is zero in degrees at least two.
Faces are induced by the actual stabilizer inclusions (after
transporting an orbit representative). The splittings (TD.4)
are compatible with such inclusions and transports by uniqueness:
two resulting sections over the same open subgroup of the same
central extension must agree. Thus no coherence is discarded
in the displayed two-row model.
\end{proof}


\subsection{Root-closed stabilizers and the remaining transgression}

\begin{proposition}[Proved: root closure and maximal cut stabilizers]
\label{G:cutroots}
Let $S$ be a locally free monoid and $K=S^{\gp}$.
The submonoid $S\subseteq K$ is root-closed: if $z^m\in S$
for $z\in K$ and an integer $m\ge1$, then $z\in S$.
Every bounded Dedekind--MacNeille cut stabilizer $K_c$ in
Proposition~\ref{G:cutstabilizers} is likewise root-closed in $K$.
If $K_c\ne1$, then, for every $1\ne h\in K_c$,
\[
 K_c=C_K(h)=N_K(K_c).
\]
Thus two nontrivial cut stabilizers either are equal or have
trivial intersection. For the finite-parameter tree $T$ of
Proposition~\ref{G:prefixtree}, if $C$ is such a nontrivial
stabilizer, the fixed-point poset $T^C$ is nonempty and has
weakly contractible nerve. Every point fixed by one nonidentity
element of $C$ has stabilizer exactly $C$.
\end{proposition}
\begin{proof}
Represent $z$ by a word in a finite free chart, and enlarge the
chart until its $m$th power equals the given positive element.
In a free group, if a positive power of a reduced word is a
positive word, then the word itself is positive. To see this,
write its cyclic reduction as $vuv^{-1}$. The reduced power is
$vu^mv^{-1}$; positivity forces $v$ to be empty, and then all
letters of $u$ are positive. The identity case is immediate.
The resulting positive representative maps to $z$ in $K$,
proving root closure of $S$.

Write a bounded cut as $(D,U)$, and suppose $g^m$ fixes it.
Choose $a\in D,b\in U$. By the cut dichotomy in
Proposition~\ref{G:cutstabilizers}, either $D$ is upward
directed or $U$ is downward directed. In the first case,
for any $d\in D$ choose $x\in D$ above $a,g^{-m}a,d$.
Then $x,g^m x$ lie in the chain $[a,b]$, so
$(x^{-1}gx)^m$ or its inverse is positive. Root closure
makes $x,gx$ comparable. Applying the order-preserving
translation by $g$ repeatedly shows that $gx$ lies between
$x$ and $g^m x$, so $gx\in D$ and $gd\le gx$ lies in $D$.
Thus $gD\subseteq D$; repeating with $g^{-1}$ gives equality.
If $U$ is downward directed, for $u\in U$ choose $x\in U$
below $b,g^{-m}b,u$. The same argument puts $gx$ between
$x,g^m x$ in $U$, and then $gu\ge gx$ lies in $U$.
Again applying the inverse gives $gU=U$. In either case
$g$ fixes the cut.

Let $1\ne h\in K_c$. This stabilizer is abelian by
Proposition~\ref{G:cutstabilizers}. If $z$ centralizes $h$,
the finitely generated subgroup $\langle z,h\rangle$ of the
locally free group $K$ is free and abelian, hence cyclic.
For $z\ne1$ some positive power of $z$ is a nonzero integral
power of $h$. Root closure of $K_c$ puts $z$ in $K_c$.
Thus $K_c=C_K(h)$.

If $z$ normalizes $K_c$, then $zhz^{-1}$ commutes with $h$.
Inside the finitely generated free subgroup $\langle z,h\rangle$,
these two nontrivial elements have the same Cayley-tree axis.
The element $z$ preserves that axis. It cannot reverse the axis:
a reversal fixes a vertex or inverts an edge, neither of which
occurs for the left action of a free group on its Cayley tree.
It therefore translates the axis and commutes with $h$.
This proves $N_K(K_c)=K_c$. The intersection assertion follows
by centralizing any nontrivial element of the intersection.

For the last assertion, any point $t$ fixed by $h\ne1$ has
$K_t=C_K(h)=C$. The nerve of $T^C$ is weakly contractible by
the fixed-path proof of Proposition~\ref{G:prefixtree}, with
any point of $T^C$ as the fixed root. Explicitly, the finitely
many turns on paths between fixed vertices are fixed, as are
their medians. Every finite set of fixed vertices therefore
lies in a finite fixed-vertex tree order whose nerve contracts
by removing leaves. Convexity of the fixed set in the ambient
pretree is not asserted or needed.
\end{proof}

\begin{corollary}[Proved: collapse of the stabilizer row]
\label{G:stabilizerrow}
In the notation of Proposition~\ref{G:treedoublemodel}, let
$\mathscr C$ be the set of distinct nontrivial groups $K_c$
for $c\in T$. It is an $H=K\rtimes G$-set under conjugation.
Choose one $C$ from every $H$-orbit, and put
\[
 A_{\mathrm{cut}}=\bigoplus_{[C]\in\mathscr C/H}C,
 \qquad X=|(NS)/G|,\qquad Y=|NT/H|,\qquad\Gamma=K_G.
\]
The degree-one stabilizer row in (TD.3) has zero homology in
positive horizontal degrees, and has $H_0=A_{\mathrm{cut}}$.
Via the two projections of the double model, there are natural isomorphisms
\[
 H_i(X;\Z)\cong H_i(Y;\Z)\quad(i\ge3)
\]
and a natural exact sequence
\[
 0\longrightarrow H_2(X;\Z)\longrightarrow H_2(Y;\Z)
 \stackrel{d_2}{\longrightarrow} A_{\mathrm{cut}}
 \longrightarrow\Gamma_{\ab}
 \longrightarrow H_1(Y;\Z)\longrightarrow0.
 \tag{TR.1}
\]
The map from each summand $C$ to $\Gamma_{\ab}$ is the
inclusion into $K$, followed by the group-coinvariant map and
abelianization. Neither $H_i(Y)$ nor the transgression $d_2$
is declared zero. In particular the remaining general
$H_2(X)$ question is exactly the kernel of the displayed
transgression.
\end{corollary}
\begin{proof}
A simplex $\sigma$ of $NT$ has nontrivial $K$-stabilizer exactly
when all its vertices belong to $T^C$ for one $C\in\mathscr C$.
In that case $K_\sigma=C$, by Proposition~\ref{G:cutroots}.
Every face of such a simplex has the same nontrivial stabilizer.
Thus the nonzero stabilizer chain groups split into the
subcomplexes associated to these distinct fixed-point posets.

For a chosen $C$, its normalizer $N_H(C)$ has kernel
$N_K(C)=C$ under projection to $G$. Its image $U$ is open:
if $1\ne h\in C$, the open group $G_h$ fixes $C=C_K(h)$
and gives a standard local section. Conjugation on $C$
descends to $U$, since $C$ is abelian. Its kernel contains
the open subgroup $G_h$, which fixes the entire rank-one
group $C$ once it fixes $h$. This open normal kernel is all
of $U$, so the extension is central. Its projection is a
quotient map, as the standard local section has open image.
Proposition~\ref{G:centralextensions} therefore gives
\[
 N_H(C)=C\times\widehat U.
\]
The subgroup $C$ acts trivially on $T^C$ by definition.
The action of $\widehat U$ on that fixed-point poset is smooth,
and $\widehat U$ has a total presentation, being isomorphic to
an open subgroup of $G$. Proposition~\ref{G:cutroots} and
Lemma~\ref{G:smoothquotient} give
\[
 |NT^C/N_H(C)|=|NT^C/\widehat U|\simeq *.
\]
The stabilizer subcomplex for this orbit is precisely
\[
 C\otimes_{\Z} C_*(NT^C/\widehat U;\Z).
\]
Indeed a simplex with stabilizer $C$ can be transported to
another such simplex only by its normalizer; the normalizer
acts trivially on the coefficient $C$, and every face map
inside this subcomplex is the identity on $C$.
This is a constant-coefficient chain complex of a weakly
contractible simplicial set. Its homology is $C$ in degree
zero and vanishes in positive degrees. Direct sums over the
$H$-orbits prove the stabilizer-row assertion.

The second page of the spectral sequence now has only the
row $E^2_{p,0}=H_p(Y;\Z)$ and the single extra group
$E^2_{0,1}=A_{\mathrm{cut}}$. The only possible nonzero
remaining differential is
$d_2:E^2_{2,0}\to E^2_{0,1}$. Reading the filtration in total
degrees one and two gives (TR.1), and all higher degrees give
the stated isomorphisms. The usual nerve presentation of
$NS/G$ identifies its fundamental group with $K_G$: its
oriented generators are $G$-orbits of elements of $S$, and
its triangular relations are their products. Thus
$H_1(X;\Z)=\Gamma_{\ab}$. The indicated edge map on each
$C$ follows from its actual inclusion into the fibre model,
so this is the specified map, not merely an abstract exact
sequence with unnamed terms.
\end{proof}


\subsection{Periodic anchors force a nonzero transgression}

\begin{example}[Proved: cut groups can die in coinvariants]
\label{G:periodicanchors}
There is an explicit diagram $M:\Cdelta\to\FreeMon$ of
finite-rank free monoids with injective transition maps for which
$T_i(M_{\ab})=0$ for every $i>0$, but a nontrivial cut stabilizer
$K_c$ of its smooth limit monoid maps noninjectively to $K_G$.
For this example, in Corollary~\ref{G:stabilizerrow},
\[
 H_2(X;\Z)=0,\qquad H_2(Y;\Z)\ne0,\qquad d_2\ne0.
\]
Thus neither cut-group injectivity nor vanishing of the quotient
$Y$'s second homology is a valid general shortcut. The example
does not refute injectivity of $d_2$; here that differential is
injective.
\end{example}
\begin{proof}
Put $E(B)=\operatorname{Emb}(V,B)$ and let $S_0(B)$ be the
interval monoid on the ordered set $E(B)$. This is free on its
successive gaps; it is trivial if $|E(B)|\le1$.
Use the reduced-arrow normal form of
Example~\ref{G:fixedgroupfailures}. Let $J(B)$ be the finite
set of nonidentity elements of $S_0(B)$ whose number of arrows
in this normal form is at most $|B|$. An increasing injection
of endpoint sets preserves both reducedness and this length.
Since a strong embedding $B\to C$ has $|B|\le|C|$, these
sets form a covariant set action $J$.

Define
\[
 M(B)=S_0(B)*\FreeMon(p_s,q_s\mid s\in J(B)).
\]
It is a finite-rank free monoid. For a strong embedding
$f:B\to C$, act on the interval monoid by its endpoint map and
put, with $d=|C|-|B|$,
\[
 M(f)(p_s)=f(s)^d p_{f(s)},\qquad
 M(f)(q_s)=f(s)^d q_{f(s)}.
 \tag{PA.1}
\]
Exponents add under composition, so this defines a functor.
The maps are injective. The interval-group inclusion is a
free-factor inclusion, using vertex bases based at any common
endpoint (and is trivial if there are fewer than two source
endpoints). The anchor labels have distinct images. After this
inclusion, (PA.1) is a positive Nielsen substitution: fix the
base group and replace each indicated anchor letter by a base
word times that letter. Its inverse in the free group removes
that base word. Thus the resulting group map, and hence its
restriction to the monoid, is injective.

There is a natural integral basis change giving
\[
 M_{\ab}\cong (S_0)_{\ab}\oplus\Z[J\times\{p,q\}].
 \tag{PA.2}
\]
It sends the anchor basis vector at $s\in J(B)$ to
$[p_s]+|B|[s]$, or $[q_s]+|B|[s]$. Under (PA.1) this is
sent to the same expression at $f(s)$ and $C$.
The basis change is invertible because it only adds base
vectors to distinct anchor vectors.

The interval summand is
\[
 (S_0(B))_{\ab}
 =\ker\bigl(\Z[E(B)]\stackrel{\aug}{\longrightarrow}\Z\bigr).
\]
The augmentation cokernel is supported on stages with $E(B)$
empty and is locally zero, because every object has a common
successor with $V$. Proposition~\ref{G:moduleinjectivization}
kills all its derived colimits. The representable
$\Z[E]$ and the constant diagram have no positive derived
colimits, and their degree-zero augmentation is an isomorphism.
The two short exact sequences through the augmentation image
therefore show that $(S_0)_{\ab}$ has zero derived colimits
in every degree. The other summand of (PA.2) is free on a
set action and is acyclic in positive degrees by totality.
This proves the asserted vanishing of $T_i(M_{\ab})$.

Let $S$ be the inclusion colimit of $M$ in the limit structure,
and $S_0$ the interval monoid on its strong-singleton orbit
$\Omega$. For every $1\ne s\in S_0$, arbitrary enlargement
of its finite strong support eventually puts $s$ in $J(B)$.
The elements
\[
 P_s=s^{|B|}p_s,\qquad Q_s=s^{|B|}q_s
\]
tagged at such a support are independent of the support, by
(PA.1). They are permuted by $G$ with the label $s$.
For every $n\ge0$ there are elements $p_{s,n},q_{s,n}\in S$
with
\[
 P_s=s^n p_{s,n},\qquad Q_s=s^n q_{s,n}:
\]
choose a support of size $N\ge n$ and use
$p_{s,n}=s^{N-n}p_s$, and likewise for $q$.
This description also shows that $S$ is the directed union of
free interval charts with finitely many independent anchor
letters and the substitutions $p_{s,n}=s p_{s,n+1}$.
In particular its action is smooth and it is locally free.

Fix a nonidentity interval $s=[x,y]$ with $x<y$ in $\Omega$.
The common lower set of $P_s,Q_s$ in the full prefix poset of
$K=S^{\gp}$ is exactly
\[
 D=\bigcup_{n\ge0}\mathord\downarrow s^n.
 \tag{PA.3}
\]
The inclusion from right to left follows from the displayed
factorizations. For the converse, a common lower vertex and
its two positive increments can be represented in a single
later free chart $M(B)$. There the two words are
$s^{|B|}p_s$ and $s^{|B|}q_s$, whose greatest common lower
vertex in the oriented Cayley tree is $s^{|B|}$, because the
terminal anchor letters are distinct. This proves (PA.3).
The associated meet cut $c$ is bounded by $1$ and $P_s$,
is defined by the two original vertices $P_s,Q_s$, and is
nonprincipal. Its lower set is preserved by $s$, since
$\bigcup_{n\ge0}\downarrow s^{n+1}=D$. Thus $1\ne s\in K_c$.
Nontriviality follows from the embedding $S\subseteq K$ and
the nonempty interval word. The absence of a largest point
in (PA.3) proves nonprincipality.

Finally, every nonidentity interval has the same image $e$
in the ordinary monoid coinvariants of $S_0$, and $e^2=e$.
Indeed, given endpoints $x<y$, freely adjoin an isolated strong
singleton $z$ between them. Both pairs $(x,z)$ and $(z,y)$
are strong free pairs and have one fixed orbit type. Writing
$e$ for the image of that type, the interval relation gives
$[x,y]=e^2$. Applying this to a strong free pair gives
$e=e^2$. Hence every such interval dies in the group
coinvariants $K_G$. In particular $s$ gives a nonzero element
of the summand $K_c$ of $A_{\mathrm{cut}}$ that maps to zero
in $\Gamma_{\ab}$.

The degreewise comparison of Proposition~\ref{G:groupcriterion}
identifies the diagram's homotopy colimit with $X=|(NS)/G|$.
The free-monoid homology formula then gives $H_2(X)=T_1(M_{\ab})=0$.
Exactness of (TR.1) makes $d_2$ injective and puts the preceding
nonzero element of $A_{\mathrm{cut}}$ in its image. Therefore
$H_2(Y)\ne0$ and $d_2\ne0$, as claimed. The exact finite
subdivision, Nielsen and basis identities are checked in
\texttt{code/G34\_periodic\_anchors.py}.
\end{proof}


\subsection{Finite-support cuts and cross-orbit comparisons}

\begin{lemma}[Proved: smooth cuts in a support lexicographic order]
\label{G:smooth-lex-cuts}
Let $G=\Aut(M_0)$ for $\Cdelta$. Fix a strong $E\subseteq F$,
an embedding $\eta:E\to M_0$, a strong flag
\[
 E=B_0\subsetneq B_1\subsetneq\cdots\subsetneq B_d=F,
 \qquad B_j=\operatorname{cl}_F(B_{j-1}\cup\{a_j\}),
\]
and signs $\epsilon_j\in\{+1,-1\}$. Give the set $\Omega$ of
strong embeddings $f:F\to M_0$ extending $\eta$ its signed
lexicographic order from this flag. If $D\subset M_0$ is finite
strong and contains $\eta E$, every $G_D$-invariant initial
segment of $\Omega$ is a union of cells in the following finite
convex partition:
\begin{enumerate}
\item for $0\le j<d$, each strong $h:B_j\to D$ extending
$\eta$, and each raw-order gap $I$ of $D$, the nonempty set
\[
 \{f\in\Omega:f|_{B_j}=h,\ f(a_{j+1})\in I\};
 \tag{OC.1}
\]
\item the singleton sets consisting of strong embeddings
$F\to D$ extending $\eta$.
\end{enumerate}
In particular membership is expressible using only equalities
between the displayed support coordinates and strict raw-order
comparisons with the finitely many points of $D$.
\end{lemma}
\begin{proof}
For any $f\in\Omega$ whose image is not contained in $D$, take
the first $j+1$ for which $fB_{j+1}\not\subseteq D$.
Then $f(a_{j+1})\notin D$. Indeed, otherwise the strong set $D$
would contain the strong closure of $fB_j\cup\{f(a_{j+1})\}$,
which is $fB_{j+1}$. Thus $f$ belongs to precisely one cell
(OC.1). Embeddings with full image in $D$ give the other cells.
These form a finite partition. Each cell is convex: fixing the
prefix $B_j$ fixes the first $j$ lexicographic coordinates,
and a raw gap, also in its reversed order, is a convex interval
for the next coordinate.

We show that a $G_D$-orbit through any point $f$ of a cell is
both cofinal and coinitial in that entire cell. Let $f'$ be
another point of the cell. Choose finite strong sets $T,Z$
containing respectively $D\cup fF$ and $D\cup f'F$. Freely
amalgamate a copy of $T$ with $Z$ over $D$. In the raw gap
containing the distinguished pivot, place the copied pivot
strictly below $f'(a_{j+1})$, or strictly above it, at will.
All remaining copied points can be placed in their prescribed
order in the gaps of $D$; this imposes no restriction on either
choice of the pivot. The resulting ordered free amalgam is in
$\Cdelta$ and embeds in $M_0$ over $Z$. Strong homogeneity
extends the copy of $T$ to an automorphism fixing $D$. We have
therefore found orbit points on both sides of $f'$ in the
signed lexicographic order. If an invariant initial segment
contains $f$, an orbit point above $f'$ shows that it contains
$f'$. If it omits $f$, an orbit point below $f'$ shows that it
omits $f'$. Membership is constant on the cell, as required.
\end{proof}

\begin{proposition}[Proved: comparisons between finitely many smooth orbits use only order]
\label{G:order-only-tree-comparisons}
Let $L$ be a smooth $G$-set with a $G$-invariant partial order
whose principal downsets are chains. Identify each of finitely
many chosen orbits with $\operatorname{Emb}(F_i,M_0)$ using its
unique finite strong support. On this finite union of orbits,
every comparison $x\preceq y$ is given by a finite Boolean
combination of equalities and raw-order comparisons between
the coordinates of the two support embeddings. No relation on
the union of the supports other than raw order and coordinate
equality is needed in these formulas.

More precisely, let $E_i\subseteq F_i$ be the component base
of the $i$th orbit from Theorem~\ref{G:alltreeorders}. If a
point of orbit $i$ lies below a point of orbit $j$, its component
label is obtained from the latter's component label by a fixed
strong embedding $E_i\to E_j$. If this embedding is an
isomorphism, each matched pair of orbit components has a
linearly ordered union. Consequently a matched pair of
components linked by a comparison has a chain union unless
the indicated inclusion of component bases is strict. No
assertion is made about incomparable components of the same
base type.

This applies, in particular, to the rooted tree order on the
completion $T$ of Proposition~\ref{G:prefixtree}, rooted at
the $G$-fixed original vertex $1$.
\end{proposition}
\begin{proof}
Fix a source orbit $\Omega_F$. Theorem~\ref{G:alltreeorders}
gives a strong $E\subseteq F$ such that its comparability
components are precisely the fibers of restriction to $E$,
and gives a signed lexicographic order in each fiber.
Fix a point $y$ of a target orbit and write its stabilizer as
$G_D$ with $D$ finite strong; the empty support is allowed.

If no point of $\Omega_F$ is below $y$, the comparison relation
on this pair of orbits is empty by transitivity on the target
orbit. Otherwise all points of $\Omega_F$ below $y$ belong to
one comparability component: any two have upper bound $y$ and
hence are comparable. Its component label $\eta:E\to M_0$
is fixed by $G_D$. Since $\operatorname{Fix}(G_D)=D$, its image
is contained in $D$. Within that fiber, the points below $y$
form a $G_D$-invariant initial segment. The preceding lemma
expresses it as a finite union of cells using only coordinate
equalities and comparisons with $D$.

For completeness, if the target orbit is identified with
$\operatorname{Emb}(D,M_0)$, the cell attached to
$h:B_j\to D$ and the gap between consecutive $d_r,d_{r+1}\in D$
transports to the formula
\[
 f|_{B_j}=g h,\qquad
 g(d_r)<f(a_{j+1})<g(d_{r+1}),
\]
where $g:D\to M_0$ represents the target point; at either end
omit the corresponding inequality. A singleton cell transports
to $f=gh$. Homogeneity and invariance make these the same
formulas for every target point, independently of the chosen
extending automorphism. A source orbit with empty support is
a singleton, and its comparison with any target orbit is
constant, which is again expressible by such a formula.

We prove the component-base assertion. If orbit $i$ has a
predecessor of one point of orbit $j$, transitivity gives a
nonempty predecessor set at every point of orbit $j$. For
comparable $y,z$ in one $j$-component, these predecessor sets
are nested, and their unique $i$-component labels are therefore
equal. Thus the label is constant on each $j$-component.
The stabilizer of that component is the pointwise stabilizer
of its $E_j$-label, so it fixes the selected $E_i$-label.
The fixed-point property puts its image inside that $E_j$;
the resulting map $E_i\to E_j$ is strong because both images
are strong in $M_0$. Equivariance makes this one fixed
embedding between the representative supports.

When this embedding is an isomorphism, the two matched
components are transitive chains under the same base
stabilizer. Every point $x$ of the $i$-component lies below
some point $z$ of the $j$-component, by transitivity applied
to a known comparison. For an arbitrary $y$ in that
$j$-component, either $z\preceq y$ and hence $x\preceq y$,
or both $x,y$ lie below $z$ and are comparable. This proves
that the union is a chain. A comparison with equal finite
base size has an isomorphism as its base embedding, proving
the stated pairwise formulation. Incomparable components can
branch even when both bases are empty: two disjoint invariant
chains above a fixed root already show this. We do not merge
components merely because a chain of other orbit types links
them.

In the final assertion, the rooted order is $x\preceq_1 y$
when $x$ belongs to the canonical tree segment from $1$ to $y$.
Its principal downsets are chains, it is $G$-invariant, and the
completion action is smooth. The proposition therefore applies.
\end{proof}

\paragraph{Scope.}
The formulas are asserted on the actual strong-embedding
orbits. Arbitrary tuples of points may have overlap patterns
which cannot occur as finite strong supports. Extending these
formulas to an unrestricted tuple template, proving that such
an extension remains a tree order, and constructing a
product-compatible exact cover are not consequences of the
proposition. The general kernel and comparison questions in
Part G remain undecided.


\subsection{Inheritance by action components}

\begin{proposition}[Proved: free-monoid tests pass to action components]
\label{G:actioninheritance}
Let $A$ be a nonempty connected small category in $\mathcal F$,
let $L:A\to\Set$ be any covariant set action, and let $B$ be
a nonempty connected component of $\operatorname{El}_A(L)$.
Then $B\in\mathcal F$. No finiteness of hom-sets or of the
sets $L(a)$ is required. More generally, if $u:A\to D$ is
an ordinary final functor, then $D\in\mathcal F$.

There is a separate degreewise assertion, not requiring
$A\in\mathcal F$: fix $i>0$ and a commutative ring $R$.
If
\[
 L_i\colim_A\bigl(R\otimes_{\Z}M_{\ab}\bigr)=0
 \quad\text{for every }M:A\to\FreeMon,
\]
then the same assertion holds with $B$ in place of $A$.
The zero ring is allowed.
\end{proposition}
\begin{proof}
The elements of $L$ belonging to $B$ form a subfunctor:
a morphism in its category of elements cannot leave its
connected component. We replace $L$ by this subfunctor,
so that $\operatorname{El}_A(L)=B$; some fibers may be empty.
Write $p:B\to A$ for the projection. Given
$N:B\to\FreeMon$, define the induced free-monoid diagram
\[
 M(a)=\mathop{*}_{l\in L(a)}N(a,l).
 \tag{AI.1}
\]
For $f:a\to a'$, its restriction to the $l$th factor is the
map $N(f):N(a,l)\to N(a',L(f)l)$ followed by the inclusion
of that target factor. This gives a functor even when
$L(f)$ identifies several indices. An empty free product is
the trivial monoid.

The universal property of free products identifies cocones
from $M$ to a fixed monoid with cocones from $N$ to that
monoid. Consequently there is a specified natural isomorphism
\[
 \colim_A^{\Mon}M\cong\colim_B^{\Mon}N.
 \tag{AI.2}
\]
At each $a$ the natural map of pointed simplicial sets
\[
 W(a):=\bigvee_{l\in L(a)}\mathrm B N(a,l)
       \longrightarrow \mathrm B M(a)
 \tag{AI.3}
\]
is a weak equivalence. Indeed each factor is the
classifying space of a free monoid, hence has the homotopy
type of the graph of its basis circles. The wedge of those
graphs is the graph of the union basis for the free product.
The indicated map induces that basis identification on
fundamental groups, and both sides are aspherical.
For arbitrary ranks and index sets this follows by the
filtered union of finite collections of basis circles.
The map (AI.3) itself is natural under all the substitutions
in (AI.1); no natural choice of the graph retractions is
needed.

The pointed bar construction gives a natural homotopy
pushout square
\[
\begin{tikzcd}
 \mathrm B B \arrow[r] \arrow[d,"\mathrm Bp"'] &
 \hocolim_{B}\mathrm B N \arrow[d] \\
 \mathrm B A \arrow[r] & \hocolim_A W.
\end{tikzcd}
 \tag{AI.4}
\]
To verify it, before identifying basepoints use the
unpointed coproduct
$\coprod_{l\in L(a)}\mathrm B N(a,l)$.
Its simplicial replacement is the simplicial replacement
over $B$: a chain in $B$ is precisely a chain in $A$ with
its starting element of $L$ specified. The coproduct of
basepoints has homotopy colimit $\mathrm B B$, and the
single common basepoint diagram has homotopy colimit
$\mathrm B A$. The inclusions of these basepoints are
simplicial cofibrations. Their objectwise pushout is $W$,
so taking the simplicial replacement and realization gives
(AI.4) as a homotopy pushout.

The inherited inclusion $\mathcal F\subseteq\{\text{total}\}$
makes $A$ total. Thus both $\mathrm B A$ and $\mathrm B B$
are weakly contractible. The right vertical map of (AI.4)
is therefore a weak equivalence. Composing it with the
homotopy colimit of (AI.3) gives a specified weak equivalence
\[
 \hocolim_B\mathrm B N\longrightarrow
 \hocolim_A\mathrm B M.
\]
Under (AI.2), the two canonical comparison maps commute
with this map. The comparison for $M$ is a weak equivalence
by $A\in\mathcal F$, so the comparison for $N$ is one by
two-out-of-three. This proves $B\in\mathcal F$.

For the final-functor assertion, each comma $d/u$ is the
category of elements of the action $a\mapsto D(d,u(a))$.
Ordinary finality makes this comma nonempty and connected;
totality of $A$ therefore makes it weakly contractible.
Thus $u$ is homotopy final as well as ordinary final. For
any free-monoid diagram on $D$, restriction along $u$
preserves both its ordinary monoid colimit and its homotopy
colimit of classifying spaces, by their specified comparison
maps. Its comparison is consequently a weak equivalence by
the comparison property on $A$.

For the degreewise assertion, (AI.1) gives the natural
isomorphism
\[
 R\otimes M_{\ab}(a)
 \cong\bigoplus_{l\in L(a)} R\otimes N_{\ab}(a,l).
\]
The ordinary bar chain complexes computing the two derived
colimits are isomorphic: in each degree the summands are
indexed by a chain in $A$ and its starting element of $L$,
equivalently by a chain in $B$, with the same coefficient
and the same faces. Thus, in every degree $i\ge0$,
\[
 L_i\colim_A(R\otimes M_{\ab})
 \cong L_i\colim_B(R\otimes N_{\ab}).
\]
This proof of the degreewise implication does not use
contractibility or totality, and also covers $R=0$.
\end{proof}

\paragraph{Group interpretation.}
In setting (S), together with the presentation invariance of
Proposition~\ref{G:groupcriterion} and the open-subgroup
presentation of Proposition~\ref{A:open-subgroups}, this
shows that the free-monoid comparison property passes to
open subgroups. It also passes along continuous dense-image
homomorphisms between groups with presentations in (S).
For the latter assertion, restrict a smooth locally free
monoid along the homomorphism. Its finite-tuple orbits are
unchanged by density, as are its monoid coinvariants. Thus
the simplicial quotient and the target in
Proposition~\ref{G:groupcriterion} are literally the same.
These are inheritance properties of $\mathcal F$, not a
determination of either inclusion in
$\{\text{exact}\}\subseteq\mathcal F\subseteq\{\text{total}\}$.


\subsection{Positive matrices do not remove the obstruction}

\begin{proposition}[Proved: positive integral stabilization]
\label{G:positive-matrix-stabilization}
Let $J$ be a small category with finite hom-sets and a degree
function $d:\operatorname{Ob}J\to\N$ such that every
nonidentity arrow strictly increases degree and there are
only finitely many objects of each bounded degree. Let
$N:J\to\Q\text{-}\mathrm{Mod}$ have finite-dimensional values.
There is a diagram $L:J\to\Ab$ of finite-rank free abelian
groups, with a specified basis at every object, such that
\[
 L\otimes_{\Z}\Q\cong\underline\Q\oplus N
 \tag{PM.1}
\]
naturally, and every nonidentity transition matrix of $L$ has
strictly positive integer entries. If all transitions of $N$
are injective, those of $L$ can also be chosen injective.

In particular there is such an injective strictly-positive-matrix
diagram on $\Cdelta$ for which
\[
 T_1(L)\otimes\Q\cong\Q.
\]
This is a statement about positive integer matrices. It does
not assert that they can be realized by a functor to free
monoids.
\end{proposition}
\begin{proof}
Put $V_a=\Q e_a\oplus N(a)$. We will choose positive rational
numbers $\lambda_a$ and lattices $L(a)\subset V_a$ so that
\[
 V(f)(t e_a+v)=
       (\lambda_b/\lambda_a)t e_b+N(f)v
       \quad(f:a\to b)
 \tag{PM.2}
\]
sends each lattice basis to a strictly positive integer
combination of the target lattice basis. Formula (PM.2) is
functorial for every choice of the $\lambda_a$, and its first
summand is naturally isomorphic to the constant rational
diagram by $t e_a\mapsto t/\lambda_a$.

Proceed by degree. At $b$, choose a rational basis
$x_1,\ldots,x_n$ of $N(b)$, and use the rational basis of $V_b$
\[
 b_0=e_b-\sum_{i=1}^n x_i,\qquad b_i=e_b+x_i
       \quad(1\le i\le n).
\]
The positive cone generated by these $n+1$ vectors has $e_b$
in its interior. More explicitly, the coefficients of
$t e_b+\sum_i v_i x_i$ in this basis are
\[
 z_0=\frac{t-\sum_i v_i}{n+1},\qquad z_i=v_i+z_0.
 \tag{PM.3}
\]
For fixed $v_i$ they are all strictly positive when $t$ is
sufficiently large.

Every already chosen source lattice basis vector has a
strictly positive $e_a$-coordinate. There are only finitely
many nonidentity arrows $a\to b$ and finitely many source
basis vectors to consider. Choose $\lambda_b>0$ sufficiently
large that (PM.3) is strictly positive for the image under
(PM.2) of every one of these vectors. All coordinates are
rational. Choose a positive integer $D_b$ clearing their
finitely many denominators, and let the lattice basis at $b$
be $b_0/D_b,\ldots,b_n/D_b$. Its image coordinates are now
strictly positive integers. Its own $e_b$-coordinates remain
positive, so the induction continues. If $n=0$, the same
argument uses the one-element basis $e_b$.

These choices yield an integral subdiagram of the rational
functor (PM.2). Compositions hold because they already hold
in $V$; the choices of positive bases impose no additional
relations. Rationalizing gives (PM.1). If $N(f)$ is injective,
(PM.2) is injective, hence so is its lattice restriction.

For the final assertion, take the rational cokernel diagram
$Q=\operatorname{coker}(P_W^{\Q}\to(P_V^{\Q})^2)$ of
Theorem~\ref{D:nonexact} for $\Cdelta$. Its values are finite
dimensional and $T_1(Q)=\Q$. Replace it by its injectivization
$Q^{\mathrm{inj}}$. Proposition~\ref{G:moduleinjectivization}
applies because $\Cdelta$ is total and preserves every derived
colimit. The quotient values remain finite dimensional, and
all transition maps are injective. A skeleton of $\Cdelta$
meets the degree hypotheses with $d(B)=|B|$: it has finitely
many isomorphism types of each bounded size, finite hom-sets,
and rigidity makes every equal-size endomorphism the identity.
Apply the construction to $N=Q^{\mathrm{inj}}$.

Rational tensor product commutes with the bar complex and
its homology. The constant rational diagram has zero positive
derived colimits because $\Cdelta$ is total. Thus (PM.1)
gives $T_1(L)\otimes\Q=T_1(Q^{\mathrm{inj}})=\Q$.
\end{proof}

\paragraph{Scope.}
Choosing the sorted word with a prescribed nonnegative
incidence column does not in general preserve composition:
substituted column blocks can interleave in the target order.
The proposition therefore isolates the missing word-level
coherence condition. It does not construct a free-monoid
counterexample, and it does not change the status of either
universal question in Part G.


\subsection{Compatible positive cones do not ensure word amalgamation}

\begin{example}[Refuted: cone-compatible bifix free-factor amalgamation]
\label{G:bifix-freefactor-AP}
The amalgamation property fails for finite-rank free monoids
with injective substitutions whose basis images are bifix
codes, even when both induced group maps are isomorphisms
and their positive abelianized matrices are identical.
Put
\[
 H=\FreeMon(x,y,z),\quad P=\FreeMon(a,b,c),\quad
 Q=\FreeMon(A,B,C),
\]
and define $i:H\to P$, $j:H\to Q$ by
\[
\begin{array}{c|ccc}
 &x&y&z\\\hline
 i&bacb&bcacb&ccacb\\
 j&BABC&BCABC&CCABC.
\end{array}
 \tag{BF.1}
\]
Both maps are injective with bifix basis images, both group
maps are isomorphisms, and their abelianized matrices are
\[
 \begin{pmatrix}1&1&1\\2&2&1\\1&2&3\end{pmatrix},
 \qquad \det=1.
 \tag{BF.2}
\]
Nevertheless there are no injective monoid maps $u:P\to S$,
$v:Q\to S$ into any locally free monoid $S$ with $ui=vj$.
\end{example}
\begin{proof}
No word in either image triple is a proper prefix or suffix
of another, so the substitutions are injective. Factor both
maps through the intermediate free monoid on $t,b,c$ by
\[
 \sigma(x)=bt,\qquad\sigma(y)=bct,\qquad\sigma(z)=cct.
\]
For $i$ send $t$ to $acb$ and leave $b,c$ unchanged; for
$j$ send $t$ to $ABC$, $b$ to $B$, and $c$ to $C$.
The bifix group basis underlying $\sigma$ is the reversal of
$\{ab,acb,acc\}$ from
\citet[Example~5.18]{GBerthe2015Bifix}. We check its inverse
explicitly. In the free group on $x,y,z$, put
\[
 c_*=zx^{-1}yz^{-1},\qquad
 t_*=c_*^{-2}z,\qquad b_*=xt_*^{-1}.
 \tag{BF.3}
\]
Their images under $\sigma^{\gp}$ are $c,t,b$, respectively.
Conversely $b_*t_*=x$, $b_*c_*t_*=y$, and
$c_*^2t_*=z$. For the middle identity, the formula for $c_*$
gives $x t_*^{-1}c_*t_*=y$ after using $z=c_*^2t_*$.
Thus $\sigma^{\gp}$ is an isomorphism. The subsequent
substitutions are Nielsen isomorphisms: their inverses send
$a$ to $t b^{-1}c^{-1}$, or $A$ to $t c^{-1}b^{-1}$.
This proves the two group-map claims; the matrix calculation
in (BF.2) is independent and immediate.

Suppose an injective amalgam existed. Passing to the group
completion of $S$ and applying the common source words
$b_*,c_*,t_*$ gives
\[
 u(b)=v(B),\qquad u(c)=v(C),\qquad
 u(a)u(c)u(b)=v(A)v(B)v(C).
 \tag{BF.4}
\]
A locally free monoid embeds in its group completion, as
proved in Proposition~\ref{G:cutstabilizers}, so these are
also equalities in $S$. Choose one finite free submonoid
containing the six target letter images. Abbreviate the
common images of $b,B$ and $c,C$ by $\beta,\gamma$ and
write $\alpha=u(a)$, $\alpha'=v(A)$. Then (BF.4) says
\[
 \alpha\gamma\beta=\alpha'\beta\gamma.
\]
Taking word lengths in this free chart gives
$|\alpha|=|\alpha'|$. Equal-length prefixes of the same
word are equal, so $\alpha=\alpha'$. Cancelling this prefix
gives $\gamma\beta=\beta\gamma$, contradicting injectivity
of $u$ because $cb\ne bc$ in $P$.

The exact script \texttt{code/G39\_bifix\_freefactor\_AP.py}
checks both bifix conditions, all eighteen inverse identities
for the intermediate and two final group maps, the identical
matrices and determinant, and the recovered words in (BF.4).
Unlike an obstruction from incompatible cones, the two
abelian group identifications here identify the entire
standard positive cones of the targets. The contradiction
uses the order of letters in the words.
\end{proof}

\paragraph{Scope.}
The example excludes the proposed blanket amalgamation
route even after imposing free-factor group maps and
compatibility of the abelian positive cones. It does not
exclude every more restrictive amalgamation class, and does
not itself give a free-monoid diagram with nonzero derived
colimit. The universal questions of Part G remain open.



\subsection{Bounded powers and the full smooth cut completion}

\begin{lemma}[Proved: bounded powers have a finite cut witness]
\label{G:bounded-power-cut}
Let $S$ be a locally free monoid, $K=S^{\gp}$, and give $K$
the prefix order. Suppose $1\ne g\in S$ and $y\in K$ satisfy
$g^n\le y$ for every integer $n\ge0$. In the bounded
Dedekind--MacNeille completion,
\[
 \sup_{n\ge0} g^n=y\wedge gy.
\]
The vertices $y$ and $gy$ are incomparable, and their displayed
meet is nonprincipal and is fixed by $g$.

Consequently every nontrivial stabilizer of an arbitrary bounded
Dedekind--MacNeille cut occurs as the stabilizer of a cut obtained
by one bounded meet or one bounded join of two original vertices,
up to left translation. The original cut need not itself have
finite parameters. If $S$ carries a smooth group action, each of
these finite witnesses is smooth.
\end{lemma}
\begin{proof}
We first exclude a positive conjugacy with bounded powers.
If $y\le gy$, write $gy=yw$ with $w\in S$. Since $1\le y$,
the elements $g,y,w$ belong to $S$. For every $n\ge0$, put
$x_n=g^{-n}y\in S$. The conjugacy equation gives
\[
 g^n x_n=y=x_n w^n.
\]
Let $F$ be the intrinsic finite-rank free hull of $g,y,w$ in
$S$. The stability rule in Lemma~\ref{G:hullfunctorial}, applied
with $p=g^n$, $q=w^n$, and the middle factor $x_n$, puts $x_n$
in $F$. Word length in this one free monoid now gives
$|y|_F\ge n|g|_F$ for all $n$, which is impossible because
$g\ne1$. The opposite comparison $gy\le y$ would give
$y=gyt$ for $t\in S$, which already contradicts length in a
single free chart containing $g,y,t$. Thus $y,gy$ are
incomparable.

The powers of $g$ have a supremum cut $c_0$, since they have
the lower bound $1$ and upper bound $y$. Their shift by $g$
is cofinal in the same set, so $gc_0=c_0$. In particular
$c_0$ is not principal, because left translation by $g\ne1$
fixes no original group vertex. Every power lies below both
$y$ and $gy$, so $c_0\le m:=y\wedge gy$.

Suppose that this inequality were strict. In terms of closed
lower sets choose $d\in D_m\setminus D_{c_0}$. The
interpolation property of Proposition~\ref{G:cutstabilizers}
gives an original vertex $k$ with
\[
 1,d\le k\le y,gy.
\]
Then $k\notin D_{c_0}$, since otherwise its lower vertex $d$
would lie there. Both $k$ and every $g^n$ belong to the chain
$[1,y]$. The possibility $k\le g^n$ would again put $k$ in
$D_{c_0}$, so $g^n\le k$ for every $n$. Also
$1\le g^{-1}k\le y$, using respectively $g\le k$ and
$k\le gy$. Thus $k,g^{-1}k$ are comparable in $[1,y]$.
After multiplying by $g$, one possibility is $gk\le k$,
excluded by positive free-chart length. The other is
$k\le gk$. Apply the preceding positive-conjugacy argument
with $k$ in place of $y$, since all powers of $g$ lie below
$k$. This is a contradiction. Hence $m=c_0$.

Now let $c=(D,U)$ be any bounded cut with nontrivial stabilizer
$C=K_c$, and choose $1\ne h\in C$. If $D$ is upward directed,
the proof of Proposition~\ref{G:cutstabilizers} supplies an
$x\in D$ for which $x,hx$ are comparable. Replace $h$ by its
inverse if necessary so that $x<hx$. Then
$g=x^{-1}hx\in S\setminus\{1\}$, and the cut $x^{-1}c$
is fixed by $g$ and has $1$ in its lower set. All powers
$g^n$ belong to that lower set and are bounded above by any
of its original upper vertices $y$. The finite meet $y\wedge
gy$ is fixed by $g$. Its stabilizer is $C_K(g)$ by
Proposition~\ref{G:cutroots}. Translating this meet by $x$
therefore gives a finite cut with stabilizer
$C_K(h)=C$.

If $D$ is not directed, then $U$ is downward directed by
Proposition~\ref{G:cutstabilizers}. Apply the order-dual
argument. Formally, inversion identifies the opposite of the
prefix order of $S$ with the prefix order of the opposite
monoid $S^{\mathrm{op}}$, which is also locally free. It
produces a bounded binary join with the required stabilizer.
Finally a cut given by two original vertices is smooth under
any smooth automorphism action, as established in
Proposition~\ref{G:cutstabilizers}.
\end{proof}

\begin{proposition}[Proved: all smooth bounded cuts give the same quotient model]
\label{G:smooth-cut-completion}
Let $G$ be the group of $\Cdelta$, let $S$ be a smooth locally
free $G$-monoid, and put $K=S^{\gp}$ and $H=K\rtimes G$.
Let $T$ be the finite-parameter tree of
Proposition~\ref{G:prefixtree}, and let $\widehat T$ be the set
of all bounded Dedekind--MacNeille cuts whose $G$-stabilizer
is open. Then $\widehat T$ is an $H$-invariant smooth median
pretree, with the same ordered-interval path construction as
$T$. Every fixed-point nerve $(N\widehat T)^U$ is weakly
contractible when $U\le G$; the analogous assertion holds for
any subgroup of $H$ having a fixed vertex. The sets of
nontrivial $K$-stabilizers occurring on $T$ and $\widehat T$
are equal.

The natural inclusion induces a weak equivalence
\[
 |NT/H|\ \longrightarrow\ |N\widehat T/H|.
\]
Consequently the double model and the transgression exact
sequence of Proposition~\ref{G:treedoublemodel} and
Corollary~\ref{G:stabilizerrow} may use all smooth bounded
cuts, with the same stabilizer summands. This does not assert
that the transgression is injective or that its source vanishes.
\end{proposition}
\begin{proof}
Every original vertex is smooth. A finite bounded meet or join
of smooth cuts is fixed by the intersection of their open
stabilizers, so $\widehat T$ is closed under these operations
and contains $T$. If $c$ is smooth and $k\in K$, then $kc$
is fixed by the intersection of $G_c$ and $G_k$, which is
open. Thus $\widehat T$ is invariant under $H$ and its
$G$-action is smooth.

The proof of Proposition~\ref{G:prefixtree} applies to this
larger set of cuts. For precision, intervals are chains by
the same interpolation argument in the original prefix poset.
Every new nonprincipal cut has an original lower and upper
vertex. Adjoining the cuts in any well-order therefore cones
off connected links and leaves the nerve connected and simply
connected. Closure under bounded binary meets and joins gives
the same associative stack reduction of finite interval paths.
It produces unique reduced paths and their medians. A finite
set of vertices together with its path turns and medians has
a finite tree hull; its nerve contracts by removing leaves.
Hence $N\widehat T$ is weakly contractible. If a group fixes
the endpoints of a path, it fixes its finitely many ordered
turns. With any fixed root, the same finite-hull contraction
proves the stated fixed-point assertion. For $U\le G$, the
original vertex $1$ supplies that root.

Lemma~\ref{G:bounded-power-cut} shows that every nontrivial
$K$-stabilizer on $\widehat T$ already occurs on $T$.
Conversely $T\subseteq\widehat T$. Denote this common set of
maximal locally cyclic subgroups by $\mathscr C$.
All the proofs of Proposition~\ref{G:treedoublemodel} and
Corollary~\ref{G:stabilizerrow} consequently apply to
$\widehat T$: stabilizers are smooth, their central extensions
split uniquely, and the fixed-point nerves needed there are
contractible. We will compare these constructions on universal
covers to obtain a weak equivalence, not only an integral
homology equivalence.

We recall the elementary quotient fundamental-group calculation.
If a group $L$ acts on a connected simply connected poset nerve
$NV$, then
\[
 \pi_1(|NV/L|)=L/\langle\!\langle L_v:v\in V\rangle\!\rangle.
\]
Here the quotient is degreewise simplicial. To verify the
formula, lift a finite quotient edge path successively, starting
at a chosen original vertex. Changes of lift are exactly
insertions of conjugates of vertex stabilizers. A path from
$v_0$ to $\ell v_0$ gives the class of $\ell$; every quotient
loop arises this way. Triangle and backtrack moves lift, after
such stabilizer changes, and simple connectedness of $NV$
kills all remaining closed lifts. This proves both surjectivity
and the displayed kernel. Poset automorphisms have no inversion
of a chain, so no further simplex-stabilizer relation is needed.

Let $J_T$ be the normal subgroup of $H$ generated by the
vertex stabilizers on $T$. It contains the standard $G$, the
stabilizer of the original vertex $1$. Modulo the normal
closure of $G$, the group $H$ is $\Gamma=K_G$.
Proposition~\ref{G:centralextensions} identifies the image
in $\Gamma$ of every affine cut stabilizer with the image of
its $K$-stabilizer. The same proof applies to smooth cuts:
the stabilizer projects onto an open subgroup, its rank-one
kernel has trivial conjugation action, and its uniquely lifted
open subgroup has trivial continuous image in the discrete
group $\Gamma$. Since the set $\mathscr C$ is unchanged, the
normal subgroup generated by all vertex stabilizers of
$\widehat T$ is exactly $J_T$. Write it as $J$, and put
$\Lambda=H/J$.

It follows that both quotient nerves have fundamental group
$\Lambda$. Moreover $\Lambda$ acts freely on the vertices
of $NT/J$ and $N\widehat T/J$: an element carrying a vertex
to its $J$-orbit differs from a $J$-element by a vertex
stabilizer, which already lies in $J$. The resulting maps to
the respective $H$-quotients are simplicial covering maps.
Indeed, choosing a lift of one specified simplex vertex fixes
a unique $\Lambda$-translate of any simplex lift; uniqueness
is compatible with all face maps. They are their universal
covers. The set of all vertex
stabilizers is invariant under $H$-conjugation, so it generates
$J$ already as a subgroup; the displayed fundamental-group
calculation applied to the $J$-action makes each covering
space simply connected.

For $V=T$ or $\widehat T$, consider the bisimplicial set
\[
 Z^V_{m,n}=(NP_m\times NV_n)/J,
\]
where $P$ is the original prefix poset. Every stabilizer of
an original $P$-simplex is conjugate to an open subgroup of
$G$ and lies in $J$. Its quotient of $NV$ is weakly
contractible by Lemma~\ref{G:smoothquotient}. Therefore the
first projection gives a natural weak equivalence
\[
 |Z^V|\ \longrightarrow\ |NP/J|
\]
for both $V$, just as in Proposition~\ref{G:treedoublemodel}.
In the other realization direction, a $V$-simplex stabilizer
in $J$ is its full stabilizer in $H$, because all vertex
stabilizers lie in $J$. Hence its fibre has the same
classifying space of its locally cyclic $K$-stabilizer as in
that proposition.

For each $C\in\mathscr C$, the normalizer splitting in
Corollary~\ref{G:stabilizerrow} is
$N_H(C)=C\times\widehat U$, with $\widehat U$ an open total
group lifted continuously. Both factors lie in $J$: the
first is a vertex stabilizer, and the second has trivial
image already in $\Gamma$. Thus $N_J(C)=N_H(C)$. The
degree-one row for $Z^V$ decomposes, now over
$J$-conjugacy classes of $C$, into
\[
 C\otimes_{\Z} C_*(NV^C/\widehat U;\Z).
\]
Each quotient nerve is weakly contractible by
Lemma~\ref{G:smoothquotient}, and inclusion from $T$ to
$\widehat T$ induces the identity on the resulting
degree-zero coefficient $C$. The degree-one row maps are
therefore homology isomorphisms in every degree.

The degree-zero rows are the chain complexes of $NT/J$ and
$N\widehat T/J$. The two total complexes have the same
homology through their natural projections to $NP/J$.
Comparing the two-row exact sequences, as in (TR.1), proves
that
\[
 |NT/J|\longrightarrow |N\widehat T/J|
\]
is an integral homology equivalence in every degree. These
spaces are simply connected CW spaces, so the relative
Hurewicz theorem makes it a weak equivalence. Together with
the common fundamental group $\Lambda$, this proves the
asserted weak equivalence of the original $H$-quotients.
\end{proof}
